% !Mode:: "TeX:UTF-8"
%!TEX program  = xelatex

\documentclass[bwprint,fontset=fandol]{gmcmthesis}

\hypersetup{hidelinks}

%========================================================
% 题目、摘要、关键词标题：隶书，小二
% 字体文件位置：Fonts/SIMLI.TTF
%========================================================

% 定义真正的隶书字体
\setCJKfamilyfont{mylishu}[
    Path = Fonts/
]{SIMLI.TTF}

% 将模板中的 \lishu 指向上传的真正隶书
\renewcommand{\lishu}{\CJKfamily{mylishu}}

%========================================================
% 将“关键词”三个字统一修改为：隶书、小二
%========================================================
\makeatletter

\renewcommand{\keywords}[1]{%
    \renewcommand{\gmcm@tokens@keywords}{#1}%
    \par
    \vskip1ex
    {\noindent\zihao{-2}\lishu
    \gmcm@cap@keywordsname：}%
    ~{\gmcm@tokens@keywords}%
}

\makeatother
\usepackage{graphicx}   % 必须，否则 \includegraphics 报错

% Support for the added equations, editable tables and figure placeholders.
\usepackage{amsmath,amssymb}
\usepackage{array,booktabs,tabularx,longtable}
\usepackage{float,placeins}

\usepackage[framemethod=TikZ]{mdframed}

%\usepackage{fontspec}  % for Consolas & Courier New
%\lstset{basicstyle=\small\fontspec{Consolas}}
%\lstset{basicstyle=\small\fontspec{Courier New}}
%\setmonofont{Consolas}
%\setcounter{tocdepth}{3}   % 调整目录深度
\usepackage{siunitx}
\usepackage{colortbl}
\definecolor{color1}{rgb}{0.78,0.88,0.99}
\definecolor{color2}{rgb}{0.36,0.62,0.84}
%\definecolor{color3}{rgb}{0.8235,0.8706,0.9373}
\definecolor{color3}{rgb}{0.88,0.92,0.96}
\definecolor{color4}{rgb}{0.96,0.97,0.98}%{0.9176,0.9373,0.9686}

% 算法
\usepackage[noend]{algpseudocode}
\usepackage{algorithmicx,algorithm}
\floatname{algorithm}{算法}
\renewcommand{\algorithmicrequire}{\textbf{输入:}}
\renewcommand{\algorithmicensure}{\textbf{输出:}}

%\numberwithin{equation}{section}
%\numberwithin{figure}{section}
%\numberwithin{table}{section}


\newcommand{\red}[1]{\textcolor{red}{#1}}
\newcommand{\blue}[1]{\textcolor{blue}{#1}}


%===================== 建模论文题目 ======================

\title{山区洪涝灾害下无人机运输与通信协同优化
}
\baominghao{\hspace{6em} 26106350114} %参赛队号
\schoolname{\hspace{5.8em} 西南大学} %学校名称
\membera{\hspace{6em} 谢文哲} %队员A
\memberb{\hspace{6em} 章林怡} %队员B
\memberc{\hspace{6em} 刘宇豪} %队员C

%=======================================================


\begin{document}

%生成标题
\maketitle

%填写摘要
\begin{abstract}

山区洪涝灾害往往伴随着\textbf{道路阻断、通信受损与能源受限}等多重约束，使传统地面运输难以及时完成应急物资配送。无人机具有机动性强、受地面道路影响小等优势，但其实际救援能力同时受到载荷、地形、能耗、配送时限、通信覆盖以及设备周转等因素制约。针对山区洪涝灾害场景下无人机应急运输与通信保障的协同优化问题，本文结合服务区空间位置、数字高程模型（DEM）、货箱需求、异构无人机参数、电池资源及无线通信链路参数，建立了由\textbf{单点运输能力分析、多机多架次调度、运输--通信联合调度和任务分区资源配置}构成的递进式优化模型，并从运输效率、任务及时性、能源消耗和资源占用等多个维度对方案进行评价。

针对问题一，为准确刻画不同机型在山区地形条件下的运输能力，首先利用\textbf{DEM栅格采样}提取调度中心与各服务区直线路径上的最高地形高程，并按照“爬升--巡航--下降”的飞行过程建立时间与能耗计算模型。在载荷影响等效航程且返航需保留安全能量余量的条件下，利用\textbf{二分搜索}求解各“机型--服务区”组合的最大安全载荷。在此基础上，将不可拆货箱的组批过程转化为同时考虑载质量、装载体积与安全能量约束的\textbf{多维装箱优化问题}，以架次数、总运输能耗和累计作业时间构造分层多目标优化模型，获得各服务区的货箱组批与机型选择方案。同时改变返航安全余量进行敏感性分析，刻画安全冗余增加对最大可运输载荷及架次数的影响规律。

针对问题二，在单架次可访问多个服务区的条件下，进一步考虑\textbf{异构无人机、共享电池、配送时限与多架次周转}之间的耦合关系。本文构建带容量、能量和时间窗约束的异构多旅行车辆路径模型，将货箱组批、服务区访问顺序、运输机型、实体无人机、电池资源及架次开始时刻进行联合决策。为降低大规模组合搜索的复杂度，采用\textbf{“可行架次生成--资源调度优化”两阶段求解框架}：首先生成满足载荷、体积和能量限制的候选运输架次，再结合医疗物资期望送达时间和首批保障截止时间，通过自适应邻域搜索与整数优化对路线及资源时序进行协调。模型能够输出逐架次运输路线、逐箱送达时刻以及无人机和电池占用情况，并通过任务时序和资源冲突检验保证方案的实际可执行性。

针对问题三，在运输调度基础上进一步引入山区地形遮挡和无线传播约束。首先依据运输无人机的连续飞行轨迹，结合DEM完成\textbf{视距（LOS）遮挡判定}，并根据双向链路预算计算运输无人机与固定网关之间的实时通信可用性。对于直连链路失效区段，在DEM覆盖区域内生成满足悬停高度约束的候选中继位置，并预计算“运输无人机--中继无人机--固定网关”两段链路的可服务关系，从而将连续空间中继部署问题转化为有限候选位置上的\textbf{时空覆盖问题}。进一步建立运输无人机与中继无人机联合调度模型，同时优化运输架次、中继悬停位置、服务时段及能源组件配置，在保证运输无人机\textbf{全飞行过程连续通信}的前提下，综合降低物资延误、联合任务完成时间、运输与中继总能耗以及设备使用强度，并通过时间离散精度的敏感性分析验证通信判定结果的稳定性。

针对问题四，为研究灾后分区独立救援条件下的资源配置问题，根据问题三中已经确定的运输架次和通信保障关系构建\textbf{服务区任务耦合图}。若多个服务区位于同一运输架次或存在不可分割的联合保障关系，则将其合并为不可拆分的任务单元，并分别建立两分区和三分区优化模型。在保持原运输与中继任务安排不变的条件下，以\textbf{资源配置规模、资源冗余、组间工作量均衡和现有库存缺口}为主要评价指标进行分区优化；同时基于任务占用时间区间计算各分区运输无人机、共享电池、中继无人机及能源组件的最大并发需求，从而得到各分区独立执行所需的最小资源配置，并识别潜在资源瓶颈及其形成原因。

综上，本文将\textbf{地形约束、运输调度、能源周转与连续通信保障}统一纳入山区应急无人机救援决策框架，在保证货箱不可拆分、任务时限、返航安全以及通信连续性的基础上，实现了由单架次能力分析到多机协同、再到运输--通信联合优化与灾区资源分区配置的逐层扩展。所建立模型具有明确的物理含义和较好的可解释性，可为山区洪涝灾害条件下无人机应急运输与通信资源协同配置提供定量化决策依据。

\keywords{山区洪涝灾害\quad  无人机应急运输\quad 异构多架次调度\quad  通信中继  \quad  协同优化}
\end{abstract}

\pagestyle{plain}

%目录 不推荐加
\maketoc

\clearpage

\section{问题重述}

\subsection{引言}   % 问题的背景
山区洪涝灾害具有突发性强、地形复杂和基础设施易受损等特点。持续强降雨可能引发山洪、滑坡及道路塌方，并进一步造成道路、供电和通信设施同时中断，使部分受灾区域在短时间内难以通过传统地面交通获得物资补给。题目所描述的广西横州市镇龙乡洪涝灾害中，就曾出现断路、断电、断网以及大量群众受困的情况，食品、饮用水、药品及抢修设备需要依靠徒步或空中方式运入灾区。 因此，利用受地面道路条件影响较小的无人机开展应急物资运输，对提高山区灾害救援效率具有重要意义。

然而，山区无人机救援并非简单的物资配送问题。一方面，复杂地形会影响无人机的飞行高度、航程及能源消耗，不同货箱的重量、体积和配送时限也使货箱组批、路径规划、机型选择和电池周转相互制约；另一方面，山体遮挡可能导致运输无人机与固定网关之间通信中断，需要利用中继无人机提供持续通信保障。由此，运输任务、通信保障和有限装备资源之间形成了明显的耦合关系。题目也明确要求同时考虑载荷、地形净空、能源储备、任务时限及通信链路等多种因素。

基于此，本文以题目给出的调度中心、15个服务区、80个不可拆货箱以及运输无人机、中继无人机、共享电池和DEM等数据为基础，围绕单点运输能力、多机多架次调度、运输与通信协同以及任务分区资源配置四个层次展开建模，在保证运输安全、配送时效和通信连续性的前提下，对有限无人机及能源资源进行合理配置，为山区洪涝灾害下的无人机应急救援提供定量化决策依据。
\subsection{问题的提出}

%\subsubsection{问题的提出内容一}

\noindent 根据题目给定的调度中心、服务区、货箱需求、DEM地形数据、运输与中继无人机参数、电池资源及通信链路参数，需要依次解决以下四个相互递进的问题。

\textbf{问题一：} 单点往返运输能力与货箱组批。针对不同机型在各服务区执行
$O_{01}\rightarrow S_i\rightarrow O_{01}$
往返任务的场景，结合DEM地形、载荷和能耗约束，确定各机型的\textbf{最大安全载荷}。在货箱不可拆分的条件下，优化\textbf{货箱组批方案}，并分析返航安全余量变化对运输能力和组批结果的影响。

\textbf{问题二：} 异构无人机多点多架次运输调度。在允许单架次访问多个服务区的条件下，综合考虑载荷、配送时限、无人机数量及共享电池周转等约束，联合优化\textbf{货箱组批、访问顺序、无人机分配和架次时序}，以提高配送及时性并降低任务完成时间、能耗和架次数。

\textbf{问题三：} 通信约束下的运输与中继联合调度。在问题二的基础上引入山区地形遮挡和通信链路约束，当运输无人机与固定网关 $G_{01}$ 无法直连时，通过中继无人机提供通信保障。联合优化\textbf{运输任务与中继无人机的悬停位置、服务时段及能源配置}，保证运输全过程通信连续，并兼顾时效与能耗。

\textbf{问题四：} 救援任务分区与资源配置。在保持问题三运输与通信任务关系不变的条件下，将 $15$ 个服务区分别划分为 $2$ 个和 $3$ 个任务组。比较不同分区方案下的\textbf{资源需求、资源冗余、工作量均衡及库存缺口}，确定更合理的分区与资源配置方案。

\clearpage

\section{模型假设与符号说明}
\subsection{模型假设}
\subsection{符号说明}
\begin{table}[htp!]
\centering
\renewcommand\arraystretch{1.2} %定义表格高度
\newcolumntype{L}{>{\quad}X}
\newcolumntype{P}[1]{>{\centering\arraybackslash}p{#1}}
\newcolumntype{C}{>{\centering \arraybackslash}X}
\newcolumntype{R}{>{\raggedright \arraybackslash}X}
\begin{tabularx}{0.9\textwidth}{|P{1.8cm}|L|}
  \hline %\toprule
  符号    &   \quad 意义 \\
  \hline %\midrule
  $ O $  &  临时调度中心    \\
  \hline
  $ S_i $  &  第 $ i $ 个服务区，$ i=1,2,\ldots,15 $    \\
  \hline
  $ g $  &  运输无人机机型，$ g \in \{A,B,C\} $    \\
  \hline
  $ m_g $  &  机型 $ g $ 的无人机空机质量/kg    \\
  \hline
  $ Q_g $  &  机型 $ g $ 的额定最大载货质量/kg    \\
  \hline
  $ V_g $  &  机型 $ g $ 的最大装载体积/m$^3$    \\
  \hline
  $ E_g $  &  机型 $ g $ 单组电池可用能量/kWh    \\
  \hline
  $ p_g $  &  机型 $ g $ 的返航安全余量比例    \\
  \hline
  $ v_g $  &  机型 $ g $ 的水平巡航速度/(m·s$^{-1}$)    \\
  \hline
  $ v_g^\uparrow, v_g^\downarrow $  &  机型 $ g $ 的爬升速度和下降速度/(m·s$^{-1}$)    \\
  \hline
  $ \eta_g $  &  爬升附加能耗效率    \\
  \hline
  $ R_g^0, R_g^F $  &  机型 $ g $ 的空载、满载标准航程/m    \\
  \hline
  $ d_{ij} $  &  节点 $ i $ 到节点 $ j $ 的水平测地距离/m    \\
  \hline
  $ H_{ij} $  &  航线 $ i \to j $ 的计划巡航海拔/m    \\
  \hline
  $ h_{ij}^\uparrow, h_{ij}^\downarrow $  &  航线 $ i \to j $ 的爬升高度、下降高度/m    \\
  \hline
  $ q $  &  当前有效载荷/kg    \\
  \hline
  $ R_g(q) $  &  载荷 $ q $ 下的等效航程/m    \\
  \hline
  $ E_{ij}(g|q) $  &  机型 $ g $ 在载荷 $ q $ 下执行航线 $ i \to j $ 的运输能耗/kWh    \\
  \hline
  $ t_{ij}(g) $  &  机型 $ g $ 执行航线 $ i \to j $ 的飞行时间/s    \\
  \hline
  $ B_i $  &  服务区 $ S_i $ 的货箱集合    \\
  \hline
  $ m_b, v_b $  &  货箱 $ b $ 的质量/kg 与体积/m$^3$    \\
  \hline
  $ q_g^{\mathrm{safe}} $  &  机型 $ g $ 服务 $ S_j $ 时的最大安全载荷/kg    \\
  \hline
  $ p $  &  某一可行班驶模式    \\
  \hline
  $ c_p $  &  班驶模式 $ p $ 对各类货箱的数量向量    \\
  \hline
  $ N $  &  全部服务区总运输频次    \\
  \hline
  $ E $  &  全部班次总运输能耗/kWh    \\
  \hline
  $ T $  &  全部班次累计作业时间/s    \\
  \hline
\end{tabularx}
\end{table}


%\begin{tabularx}{\textwidth-18pt}{XXX}
%\hline
%Input & Output& Action return \\
%\hline
%DNF &  simulation & jsp\\
%\hline
%\end{tabularx}


\section{问题一的建模与求解}

\subsection{问题一分析}
问题一要求在不考虑实体无人机和共享电池调度的条件下，研究“临时调度中心 $O01$---单个服务区 $S_i$---$O01$”的直接往返运输能力与货箱组批问题。由于每一架次只能服务一个服务区，15 个服务区之间不存在跨区组批，因此该问题可以先对每个“机型—服务区”组合建立独立的飞行时间与能耗模型，再在各服务区内部求解不可拆货箱的最优组批，最后汇总得到全局结果。
与普通装箱问题相比，本题的关键难点在于：第一，载荷上限不仅由机型额定载质量决定，还受 DEM 地形引起的爬升能耗、航程衰减和返航安全余量共同约束；第二，货箱不可拆分，同时受到质量和体积双重限制；第三，题目要求综合考虑往返架次数、运输能耗与累计作业时间，三者之间存在明显的层级权衡。因此，本文采用“物理可行性计算—可行组批模式枚举—状态压缩动态规划”的思路进行精确求解，整体流程如图~\ref{fig:problem1_flow}所示。
\begin{figure}[htbp!]
  \centering
  \includegraphics[width=0.85\textwidth]{figures/Q1_FIg1.png}
  \caption{问题一建模与求解总框架}
  \label{fig:problem1_flow}
\end{figure}
\subsection{基本约定}
（1）每一运输架次均采用 $O01 \to S_i \to O01$ 的直接往返形式，只服务一个服务区；同一服务区允许使用多个架次完成配送，但不得跨服务区组批。

（2）任意两个任务节点之间的水平航迹取两节点连线；巡航海拔由该连线经过的 $30\,\mathrm{m}$ DEM 像元最高地形高程加 $50\,\mathrm{m}$ 安全净空确定。

（3）$O01$ 的作业高度取地面海拔，服务区作业高度取地面海拔以上 $30\,\mathrm{m}$；每个航段均按"爬升—水平巡航—下降"三个阶段计算飞行时间。

（4）下降阶段不单独计入附加势能消耗；运输总能耗由载荷相关水平航程能耗与爬升附加能耗构成。

（5）问题一不考虑实体无人机数量、共享电池数量及其充电周转，只考察单架次物理可行性与货箱组批；货箱不可拆分且每个货箱只能被安排一次。
% === Problem one: imported from the reference DOCX, sections 3.3--3.8 ===
% Data and conclusions are retained from the reference; no recomputation.
% Notation follows the existing symbol table: rho_g -> p_g,
% E_ij^g(q) -> E_ij(g|q), t_ij^g -> t_ij(g).
% The safe-payload symbol includes i to distinguish service areas.
\FloatBarrier

\subsection{航段时间与能耗模型}

\subsubsection{基于 DEM 的航段净空模型}

对于节点 $i$ 到节点 $j$ 的运输航段，沿节点连线逐个检查其穿越的 DEM 像元。设这些像元的最高地形高程为 $\max(h_k)$，则计划巡航海拔为

\begin{equation}
H_{ij} = \max_{k\in L_{ij}} h_k + 50.
\label{eq:q1-clearance}
\end{equation}

其中，$L_{ij}$ 表示航段 $i\to j$ 穿越的 DEM 像元集合。设节点 $i$、$j$ 的实际作业海拔分别为 $z_i^o$、$z_j^o$，则爬升高度和下降高度分别为

\begin{equation}
h_{ij}^{\uparrow} = H_{ij}-z_i^{o},\qquad
h_{ij}^{\downarrow} = H_{ij}-z_j^{o}.
\label{eq:q1-vertical-height}
\end{equation}

该处理方式能够将山区地形起伏直接映射到无人机的爬升需求和能量消耗中，避免仅按平面距离估计运输能力。

\subsubsection{载荷相关等效航程与运输能耗}

题目给出了各机型空载标准航程与满载标准航程。为描述载荷增加导致的航程衰减，采用题目配套计算规则中的载荷相关等效航程关系：

\begin{equation}
R_g(q)=R_g^0-\bigl(R_g^0-R_g^F\bigr)
\left(\frac{q}{Q_g}\right)^{3/2},\qquad 0\le q\le Q_g.
\label{eq:q1-range}
\end{equation}

于是，机型 $g$ 携带载荷 $q$ 执行航段 $i\to j$ 时，水平巡航能耗可按“可用电池能量$\times$航段距离/等效航程”折算；同时考虑爬升势能，航段总运输能耗为

\begin{equation}
E_{ij}(g\mid q)
= E_g\frac{d_{ij}}{R_g(q)}
+\frac{(m_g+q)g_0h_{ij}^{\uparrow}}{\eta_g\,3.6\times 10^6}.
\label{eq:q1-segment-energy}
\end{equation}

式中 $g_0=9.81\,\mathrm{m}/\mathrm{s}^{2}$；第二项除以 $3.6\times10^6$ 将焦耳换算为 kWh。由于题目规定下降能耗效率取 0，故下降阶段不再单独增加附加能耗。

\subsubsection{飞行时间与单架次作业时间}

机型 $g$ 在航段 $i\to j$ 上的飞行时间由爬升、巡航和下降三个阶段构成：

\begin{equation}
t_{ij}(g)
=\frac{h_{ij}^{\uparrow}}{v_g^{\uparrow}}
+\frac{d_{ij}}{v_g}
+\frac{h_{ij}^{\downarrow}}{v_g^{\downarrow}}.
\label{eq:q1-flight-time}
\end{equation}

对于问题一中的直接往返架次，若某批次含 $n$ 个货箱，则该架次累计作业时间还包括起飞准备、逐箱装载、服务区作业以及逐箱卸载。因此

\begin{equation}
\begin{aligned}
T_p={}&t_g^{\mathrm{prep}}+n\,t_g^{\mathrm{load}}
+t_{\mathrm{O01},i}(g)+t_{i,\mathrm{O01}}(g)\\
&+t_g^{\mathrm{service}}+n\,t_g^{\mathrm{unload}}.
\end{aligned}
\label{eq:q1-operation-time}
\end{equation}

\subsection{单点往返最大安全载荷模型}

问题一首先要求计算三种机型在各服务区执行直接往返任务时的最大安全载荷。起飞时无人机携带全部货物，到达服务区后完成投送，返航阶段按空载计算。因此，机型 $g$ 向服务区 $S_i$ 运送载荷 $q$ 时的一次往返总能耗为

\begin{equation}
E_{gi}^{\mathrm{rt}}(q)
=E_{\mathrm{O01},i}(g\mid q)+E_{i,\mathrm{O01}}(g\mid 0).
\label{eq:q1-roundtrip-energy}
\end{equation}

为保证返航安全，任务结束时必须至少保留 $p_g$ 比例的可用电量，即

\begin{equation}
E_{gi}^{\mathrm{rt}}(q)\le (1-p_g)E_g.
\label{eq:q1-energy-limit}
\end{equation}

因此，机型 $g$ 在服务区 $S_i$ 的最大安全载荷定义为

\begin{equation}
q_{gi}^{\mathrm{safe}}=\max\left\{q\;\middle|\;
0\le q\le Q_g,\quad E_{gi}^{\mathrm{rt}}(q)\le (1-p_g)E_g
\right\}.
\label{eq:q1-safe-payload}
\end{equation}

若额定载荷 $Q_g$ 已满足式~\eqref{eq:q1-energy-limit}，则 $q_{gi}^{\mathrm{safe}}=Q_g$；否则在 $[0,Q_g]$ 上利用单调求根得到使式~\eqref{eq:q1-energy-limit}恰好取等号的载荷。若 $q=0$ 时仍不能满足返航余量，则判定该“机型—服务区”组合在当前安全余量下不可达。
\begin{figure}[H]
  \centering
  \includegraphics[width=0.85\textwidth]{figures/payload_distance.pdf}

  \caption{典型困难服务区 S008 的安全载荷随返航余量变化}
  \label{fig:q1-s008-sensitivity}
\end{figure}

\subsection{货箱组批优化模型}

\subsubsection{可行组批约束}

对服务区 $S_i$ 的货箱集合 $B_i$，任意候选批次 $p\subseteq B_i$ 只能选择一种机型 $g_p$。设该批次的货箱总质量和总体积分别为 $M_p$ 与 $V_p$，则必须同时满足

\begin{equation}
M_p=\sum_{b\in p}m_b\le Q_{g_p},\qquad
V_p=\sum_{b\in p}v_b\le V_{g_p}.
\label{eq:q1-capacity}
\end{equation}

此外，该批次还必须满足返航安全能量约束

\begin{equation}
E_{\mathrm{O01},i}(g_p\mid M_p)
+E_{i,\mathrm{O01}}(g_p\mid 0)
\le (1-p_{g_p})E_{g_p}.
\label{eq:q1-batch-energy}
\end{equation}

式~\eqref{eq:q1-capacity}和式~\eqref{eq:q1-batch-energy}共同决定一个货箱组合是否可以由某一机型完成。需要特别指出，最大安全载荷只约束“质量”这一连续变量，实际组批仍可能首先受到装载体积限制，因此不能简单依据载重上限进行贪心装箱。

\subsubsection{多目标优化与优先关系}

题目要求在完成全部货箱交付的前提下，综合考虑往返架次数、总运输能耗和累计作业时间。考虑到应急运输中减少架次数能够直接降低任务组织复杂度和设备切换需求，本文采用字典序优化，将架次数作为第一优先级，在最少架次方案集合内继续最小化运输能耗，最后最小化累计作业时间：

\begin{equation}
\min_{\mathrm{lex}}\;(N,E,T).
\label{eq:q1-objective}
\end{equation}

其中 $N=\sum_p x_p$，$E=\sum_p E_p x_p$，$T=\sum_p T_p x_p$；$x_p\in\{0,1\}$ 表示可行组批模式 $p$ 是否被选用。每个货箱必须且只能出现一次，即

\begin{equation}
\sum_{p:\,b\in p}x_p=1,\qquad \forall b\in B_i.
\label{eq:q1-exact-cover}
\end{equation}

\subsection{精确动态规划求解方法}

由于问题一禁止跨服务区组批，因此 15 个服务区可以分别精确求解。为降低状态维度，将同一服务区内“物资类别、单箱质量、单箱体积”完全相同的货箱归为一类。设该服务区共有 $K$ 类货箱，第 $k$ 类货箱数量为 $n_k$，定义状态向量 $\boldsymbol{x}=(x_1,x_2,\ldots,x_K)$，其中 $x_k$ 表示尚未分配的第 $k$ 类货箱数量。

首先枚举所有满足质量、体积和返航余量约束的可行组批模式 $p$，并记录其货箱数量向量 $\boldsymbol{c}_p$、机型 $g_p$、能耗 $E_p$ 和作业时间 $T_p$。令 $F(\boldsymbol{x})$ 表示完成状态 $\boldsymbol{x}$ 中全部剩余货箱所需的最优代价向量，则有递推关系

\begin{equation}
F(\boldsymbol{x})=\min_{\mathrm{lex}}\left\{
(1,E_p,T_p)+F(\boldsymbol{x}-\boldsymbol{c}_p)
\;\middle|\;
\boldsymbol{0}<\boldsymbol{c}_p\le\boldsymbol{x},\quad
p\,\text{可行}
\right\},
\label{eq:q1-dp}
\end{equation}

\begin{equation}
F(\boldsymbol{0})=(0,0,0).
\label{eq:q1-dp-initial}
\end{equation}

采用记忆化搜索保存已计算状态，能够避免对同一剩余货箱状态重复求解。各服务区求得的最优组批结果彼此独立，其架次数、能耗和作业时间可以直接相加，从而得到问题一的全局字典序最优解。为进一步验证动态规划结果，本文又以相同的可行组批模式建立 0-1 整数规划进行交叉检查，15 个服务区的最少架次数和最小能耗均与动态规划结果一致。

\subsection{计算结果与分析}

\subsubsection{三种机型的最大安全载荷}

在题目默认的 20\% 返航安全余量下，三种机型对 15 个服务区的最大安全载荷如图~\ref{fig:q1-safe-payload} 和表~\ref{tab:q1-safe-payload} 所示。A 型无人机在全部服务区均可达到其 $25\,\mathrm{kg}$ 的额定载荷；B 型除 S008 外均可达到 $30\,\mathrm{kg}$，S008 的最大安全载荷降至 $28.801\,\mathrm{kg}$；C 型受地形和航程影响最为明显，在 S002、S003、S004、S008 和 S012 的最大安全载荷分别为 68.860、68.206、63.697、58.904 和 $68.117\,\mathrm{kg}$，其余服务区仍可达到 $80\,\mathrm{kg}$ 的额定上限。

% Figure placeholder only; no external image is required.
\begin{figure}[H]
  \centering
\includegraphics[width=0.85\textwidth]{figures/3.2.pdf}

  \caption{最大安全载荷热力图}
  \label{fig:q1-safe-payload}
\end{figure}

\begin{table}[H]
  \centering
  \small
  \setlength{\tabcolsep}{5pt}
  \renewcommand{\arraystretch}{1.15}
  \caption{20\%返航安全余量下各服务区最大安全载荷（kg）}
  \label{tab:q1-safe-payload}
  \begin{tabular*}{0.75\linewidth}{@{\extracolsep{\fill}}crrr@{}}
    \toprule
    服务区 & A型 & B型 & C型 \\
    \midrule
    S001 & 25.000 & 30.000 & 80.000 \\
    S002 & 25.000 & 30.000 & 68.860 \\
    S003 & 25.000 & 30.000 & 68.206 \\
    S004 & 25.000 & 30.000 & 63.697 \\
    S005 & 25.000 & 30.000 & 80.000 \\
    S006 & 25.000 & 30.000 & 80.000 \\
    S007 & 25.000 & 30.000 & 80.000 \\
    S008 & 25.000 & 28.801 & 58.904 \\
    S009 & 25.000 & 30.000 & 80.000 \\
    S010 & 25.000 & 30.000 & 80.000 \\
    S011 & 25.000 & 30.000 & 80.000 \\
    S012 & 25.000 & 30.000 & 68.117 \\
    S013 & 25.000 & 30.000 & 80.000 \\
    S014 & 25.000 & 30.000 & 80.000 \\
    S015 & 25.000 & 30.000 & 80.000 \\
    \bottomrule
  \end{tabular*}
\end{table}

从结果可见，最大安全载荷并非仅由额定载质量决定。对 A 型而言，额定载荷较小，因此 20\% 余量下主要受到 $25\,\mathrm{kg}$ 载重上限约束；对重载 C 型而言，额定载荷虽然达到 $80\,\mathrm{kg}$，但在距离较远或需要较大爬升的服务区，载荷增加会同时降低等效航程并增加爬升势能消耗，因此安全载荷明显下降。S008 对 B、C 两型均形成约束，是后续敏感性分析中具有代表性的困难服务区。

\subsubsection{最优货箱组批方案}

在 20\% 返航安全余量和“架次数—能耗—累计作业时间”字典序目标下，80 个货箱共被划分为 18 个直接往返架次，总运输能耗为 $59.1313\,\mathrm{kWh}$，累计作业时间为 $9.104\,\mathrm{h}$，其中纯飞行时间为 $5.358\,\mathrm{h}$。最终使用 B 型 9 架次、C 型 9 架次，A 型未被采用。各服务区的组批结构如图~\ref{fig:q1-batching} 所示，完整架次结果见表~\ref{tab:q1-batching}。

% Figure placeholder only; no external image is required.
\begin{figure}[H]
  \centering
  
  \includegraphics[width=0.85\textwidth]{figures/flight_loading_red_blue.pdf}

  \caption{18 架次货箱组成图}
  \label{fig:q1-batching}
\end{figure}

\FloatBarrier
% The sortie table can span pages; column headers repeat automatically.
\begingroup
\small
\setlength{\tabcolsep}{4pt}
\renewcommand{\arraystretch}{1.2}
\setlength{\LTleft}{\fill}
\setlength{\LTright}{\fill}
\begin{longtable}{@{}lcclrrrrr@{}}
  \caption{问题一最优组批与运输结果}\label{tab:q1-batching}\\
  \toprule
    架次 & 服务区 & 机型 & 货箱组成 & \shortstack{载荷\\ /$\mathrm{kg}$} & \shortstack{体积\\ /$\mathrm{m}^{3}$} & \shortstack{能耗\\ /$\mathrm{kWh}$} & \shortstack{作业时间\\ /$\mathrm{min}$} & \shortstack{返航SOC\\ /$\%$} \\
  \midrule
  \endfirsthead
  \multicolumn{9}{c}{表~\thetable（续）}\\
  \toprule
    架次 & 服务区 & 机型 & 货箱组成 & \shortstack{载荷\\ /$\mathrm{kg}$} & \shortstack{体积\\ /$\mathrm{m}^{3}$} & \shortstack{能耗\\ /$\mathrm{kWh}$} & \shortstack{作业时间\\ /$\mathrm{min}$} & \shortstack{返航SOC\\ /$\%$} \\
  \midrule
  \endhead
  \midrule
  \multicolumn{9}{r}{\footnotesize 续下页}\\
  \endfoot
  \bottomrule
  \endlastfoot
    Q1-001 & S001 & C & \shortstack[l]{\strut $\mathrm{FOD}\times 2$\quad $\mathrm{HYG}\times 2$\\\strut $\mathrm{MED}\times 2$\quad $\mathrm{WAT}\times 3$} & 76 & 0.231 & 2.918 & 28.20 & 63.5 \\[3pt]
    Q1-002 & S001 & C & \shortstack[l]{\strut $\mathrm{FOD}\times 1$\quad $\mathrm{WAT}\times 5$} & 78 & 0.163 & 2.996 & 24.90 & 62.6 \\[3pt]
    Q1-003 & S002 & B & \shortstack[l]{\strut $\mathrm{FOD}\times 1$\quad $\mathrm{MED}\times 1$\\\strut $\mathrm{WAT}\times 1$} & 25 & 0.067 & 2.713 & 30.28 & 32.2 \\[3pt]
    Q1-004 & S002 & C & \shortstack[l]{\strut $\mathrm{FOD}\times 1$\quad $\mathrm{HYG}\times 1$\\\strut $\mathrm{WAT}\times 3$} & 56 & 0.144 & 5.721 & 34.03 & 28.5 \\[3pt]
    Q1-005 & S003 & B & \shortstack[l]{\strut $\mathrm{FOD}\times 1$\quad $\mathrm{MED}\times 1$\\\strut $\mathrm{WAT}\times 1$} & 25 & 0.067 & 2.780 & 34.68 & 30.5 \\[3pt]
    Q1-006 & S003 & C & \shortstack[l]{\strut $\mathrm{FOD}\times 1$\quad $\mathrm{HYG}\times 1$\\\strut $\mathrm{WAT}\times 3$} & 56 & 0.144 & 5.765 & 39.50 & 27.9 \\[3pt]
    Q1-007 & S004 & C & \shortstack[l]{\strut $\mathrm{FOD}\times 1$\quad $\mathrm{HYG}\times 1$\\\strut $\mathrm{MED}\times 1$\quad $\mathrm{WAT}\times 3$} & 59 & 0.156 & 6.153 & 38.35 & 23.1 \\[3pt]
    Q1-008 & S005 & C & \shortstack[l]{\strut $\mathrm{FOD}\times 1$\quad $\mathrm{HYG}\times 1$\\\strut $\mathrm{MED}\times 1$\quad $\mathrm{WAT}\times 3$} & 59 & 0.156 & 4.222 & 31.29 & 47.2 \\[3pt]
    Q1-009 & S006 & C & \shortstack[l]{\strut $\mathrm{FOD}\times 1$\quad $\mathrm{HYG}\times 1$\\\strut $\mathrm{MED}\times 1$\quad $\mathrm{WAT}\times 3$} & 59 & 0.156 & 2.598 & 26.03 & 67.5 \\[3pt]
    Q1-010 & S007 & C & \shortstack[l]{\strut $\mathrm{FOD}\times 1$\quad $\mathrm{HYG}\times 1$\\\strut $\mathrm{MED}\times 1$\quad $\mathrm{WAT}\times 2$} & 45 & 0.129 & 3.411 & 31.79 & 57.4 \\[3pt]
    Q1-011 & S008 & C & \shortstack[l]{\strut $\mathrm{FOD}\times 1$\quad $\mathrm{HYG}\times 1$\\\strut $\mathrm{MED}\times 1$\quad $\mathrm{WAT}\times 2$} & 45 & 0.129 & 5.836 & 36.66 & 27.0 \\[3pt]
    Q1-012 & S009 & B & \shortstack[l]{\strut $\mathrm{FOD}\times 1$\quad $\mathrm{MED}\times 1$\\\strut $\mathrm{WAT}\times 1$} & 25 & 0.067 & 2.096 & 25.97 & 47.6 \\[3pt]
    Q1-013 & S010 & B & \shortstack[l]{\strut $\mathrm{FOD}\times 1$\quad $\mathrm{MED}\times 1$\\\strut $\mathrm{WAT}\times 1$} & 25 & 0.067 & 1.823 & 25.92 & 54.4 \\[3pt]
    Q1-014 & S011 & B & \shortstack[l]{\strut $\mathrm{FOD}\times 1$\quad $\mathrm{MED}\times 1$\\\strut $\mathrm{WAT}\times 1$} & 25 & 0.067 & 1.074 & 19.81 & 73.2 \\[3pt]
    Q1-015 & S012 & B & \shortstack[l]{\strut $\mathrm{FOD}\times 1$\quad $\mathrm{MED}\times 1$\\\strut $\mathrm{WAT}\times 1$} & 25 & 0.067 & 2.752 & 32.05 & 31.2 \\[3pt]
    Q1-016 & S013 & B & \shortstack[l]{\strut $\mathrm{FOD}\times 1$\quad $\mathrm{MED}\times 1$\\\strut $\mathrm{WAT}\times 1$} & 25 & 0.067 & 1.825 & 25.17 & 54.4 \\[3pt]
    Q1-017 & S014 & B & \shortstack[l]{\strut $\mathrm{FOD}\times 1$\quad $\mathrm{MED}\times 1$\\\strut $\mathrm{WAT}\times 1$} & 25 & 0.067 & 2.140 & 31.16 & 46.5 \\[3pt]
    Q1-018 & S015 & B & \shortstack[l]{\strut $\mathrm{FOD}\times 1$\quad $\mathrm{MED}\times 1$\\\strut $\mathrm{WAT}\times 1$} & 25 & 0.067 & 2.307 & 30.47 & 42.3 \\[3pt]
\end{longtable}
\endgroup

注：同一服务区内同类别货箱的质量与体积参数一致，因此“货箱组成”按类别计数表示；FOD、HYG、MED、WAT 分别表示应急食品、生活卫生用品、医疗物资和饮用水。

结果显示，只有 S001、S002 和 S003 需要 2 个架次，其余 12 个服务区均可在 1 个架次内完成。其中 S001 总需求质量为 $154\,\mathrm{kg}$，必须由两架次 C 型承担；S002、S003 的总需求质量均为 $81\,\mathrm{kg}$，最优方案均采用“$25\,\mathrm{kg}$ 的 B 型架次 + $56\,\mathrm{kg}$ 的 C 型架次”。S004—S008 主要由 C 型单架次完成，S009—S015 则均由 B 型单架次完成。

值得注意的是，A 型无人机在本方案中没有被选中，这并不意味着其航程不可行，而是质量与体积约束共同作用的结果。以多个服务区中常见的“饮用水 $14\,\mathrm{kg}$ + 医疗物资 $3\,\mathrm{kg}$ + 应急食品 $8\,\mathrm{kg}$”组合为例，总质量恰为 $25\,\mathrm{kg}$，但总体积为 $0.067\,\mathrm{m}^{3}$，超过 A 型 $0.060\,\mathrm{m}^{3}$ 的装载体积，却可以装入 B 型 $0.073\,\mathrm{m}^{3}$ 的货舱。因此，若只按载重进行组批会错误高估 A 型的可用性。
\begin{figure}[H]
  \centering
  \fbox{%
    \begin{minipage}[c][4.2cm][c]{0.88\linewidth}
      \centering
      \small 图片占位框\par\medskip
      能耗优先方案相对本文方案的指标变化
    \end{minipage}%
  }
  \caption{载荷—距离—能耗响应曲面}
  \label{fig:q1-objective-comparison}
\end{figure}

\subsubsection{架次数、能耗与作业时间的权衡}

为检验字典序优先关系是否合理，进一步分别以能耗和累计作业时间作为第一目标进行对照求解，结果见表~\ref{tab:q1-objective-comparison}。

\begin{table}[H]
  \centering
  \small
  \setlength{\tabcolsep}{5pt}
  \renewcommand{\arraystretch}{1.15}
  \caption{不同优化优先级下的结果对比}
  \label{tab:q1-objective-comparison}
  \begin{tabular*}{\linewidth}{@{\extracolsep{\fill}}lrrrrr@{}}
    \toprule
    优化策略 & 架次数 & \shortstack{总能耗\\ /$\mathrm{kWh}$} & \shortstack{累计作业时间\\ /$\mathrm{h}$} & B型架次 & C型架次 \\
    \midrule
    架次数优先（本文） & 18 & 59.1313 & 9.104 & 9 & 9 \\
    能耗优先 & 19 & 59.0339 & 9.567 & 11 & 8 \\
    时间优先 & 18 & 59.1313 & 9.104 & 9 & 9 \\
    \bottomrule
  \end{tabular*}
\end{table}

% Figure placeholder only; no external image is required.
\begin{figure}[H]
  \centering
  
  \includegraphics[width=0.85\textwidth]{figures/运输与能耗变化率对比图.pdf}

  \caption{优化优先关系的代价比较}
  \label{fig:q1-objective-comparison}
\end{figure}

由表~\ref{tab:q1-objective-comparison} 可知，若将能耗置于第一优先级，总能耗可由 $59.1313\,\mathrm{kWh}$ 降至 $59.0339\,\mathrm{kWh}$，仅降低 0.165\%；但运输架次数由 18 增至 19，增加 5.56\%，累计作业时间也增加 5.08\%。而以累计作业时间为第一目标时，所得方案与本文字典序方案一致。可见，在本算例中，以减少架次数作为第一优先级只付出了极小的能耗代价，却能够降低任务批次数和总体作业时间，因而更符合灾后应急运输的组织需求。

\subsubsection{返航安全余量敏感性分析}

为分析安全裕度变化对运输能力和组批结果的影响，将返航安全余量从 10\% 逐步提高到 40\%，并在每个余量水平下重新计算安全载荷和最优组批。整体结果见表~\ref{tab:q1-reserve-sensitivity} 和图~\ref{fig:q1-reserve-sensitivity}。

\begin{table}[H]
  \centering
  \small
  \setlength{\tabcolsep}{5pt}
  \renewcommand{\arraystretch}{1.15}
  \caption{不同返航安全余量下的最优组批结果}
  \label{tab:q1-reserve-sensitivity}
  \begin{tabular*}{0.96\linewidth}{@{\extracolsep{\fill}}ccrrr@{}}
    \toprule
    返航余量 & 可行性 & 总架次数 & \shortstack{总能耗\\ /$\mathrm{kWh}$} & \shortstack{累计作业时间\\ /$\mathrm{h}$} \\
    \midrule
    10\% & 可行 & 18 & 59.131 & 9.104 \\
    15\% & 可行 & 18 & 59.131 & 9.104 \\
    20\% & 可行 & 18 & 59.131 & 9.104 \\
    25\% & 可行 & 19 & 61.073 & 9.602 \\
    30\% & 可行 & 20 & 67.228 & 10.163 \\
    35\% & 可行 & 25 & 75.608 & 12.607 \\
    40\% & 不可行 & - & - & - \\
    \bottomrule
  \end{tabular*}
\end{table}

当返航余量处于 10\%—20\% 时，虽然部分“机型—服务区”的连续安全载荷已经逐渐下降，但尚未跨越当前最优货箱组合的关键阈值，因此最优组批始终保持 18 架次，总能耗也不发生变化。当余量提高至 25\%、30\% 和 35\% 后，一些原本可以合并的货箱组合依次失去可行性，总架次数分别跃升至 19、20 和 25 架次，总能耗也由 $59.131\,\mathrm{kWh}$ 增长至 61.073、67.228 和 $75.608\,\mathrm{kWh}$。余量达到 40\% 时，全部货箱已经无法在现有三类机型下完成配送，整体问题变为不可行。

这一结果说明，由于货箱不可拆分，返航余量对组批方案的影响具有明显的“阈值效应”：安全载荷本身随余量连续下降，但架次数并不会连续变化，而是在安全载荷跌破某一离散货箱组合质量时发生跳变。

% Figure placeholder only; no external image is required.
\begin{figure}[H]
  \centering
  \includegraphics[width=0.85\textwidth]{figures/返航余量与组批结构、能耗、时间.png}
  \caption{返航余量与组批结构、能耗、时间}
  \label{fig:q1-s008-sensitivity}
\end{figure}

以 S008 为例，20\% 返航余量下 A、B、C 三型的最大安全载荷分别为 25.000、28.801 和 $58.904\,\mathrm{kg}$；余量提高到 30\% 后分别降至 14.213、21.001 和 $36.814\,\mathrm{kg}$。继续提高到 35\% 时，A 型即使空载也无法满足返航余量，而 C 型最大安全载荷仅剩 $15.683\,\mathrm{kg}$；到 40\% 时，A、C 两型均无法完成该服务区的安全往返，B 型也只能携带约 $1.397\,\mathrm{kg}$。该现象解释了高安全余量下架次数迅速增加直至整体不可行的原因。
\begin{figure}[H]
  \centering
  \includegraphics[width=0.85\textwidth]{figures/s800.png}
  \caption{S008 安全载荷敏感性}
  \label{fig:q1-s008-sensitivity}
\end{figure}

\subsubsection{求解正确性与精确性检验}

为避免仅依赖单一算法得到的结果，本文对问题一进行了两类独立核验。首先，基于 DEM 对节点间航段进行逐像元地形检查，确保巡航海拔与地形最高点的计算一致；其次，以动态规划枚举得到的可行组批模式为基础，对每个服务区单独建立 0-1 整数规划模型，并检查最少架次数与对应最小能耗。15 个服务区的整数规划最优间隙均为 0，且结果全部与动态规划一致。

\begin{table}[H]
  \centering
  \small
  \setlength{\tabcolsep}{5pt}
  \renewcommand{\arraystretch}{1.15}
  \caption{动态规划与整数规划交叉验证结果}
  \label{tab:q1-validation}
  \begin{tabular*}{0.92\linewidth}{@{\extracolsep{\fill}}crrcc@{}}
    \toprule
    服务区 & 最少架次数 & \shortstack{对应能耗\\ /$\mathrm{kWh}$} & MILP gap & 一致性 \\
    \midrule
    S001 & 2 & 5.9134 & 0.0 & 通过 \\
    S002 & 2 & 8.4346 & 0.0 & 通过 \\
    S003 & 2 & 8.5449 & 0.0 & 通过 \\
    S004 & 1 & 6.1529 & 0.0 & 通过 \\
    S005 & 1 & 4.2222 & 0.0 & 通过 \\
    S006 & 1 & 2.5980 & 0.0 & 通过 \\
    S007 & 1 & 3.4109 & 0.0 & 通过 \\
    S008 & 1 & 5.8360 & 0.0 & 通过 \\
    S009 & 1 & 2.0962 & 0.0 & 通过 \\
    S010 & 1 & 1.8232 & 0.0 & 通过 \\
    S011 & 1 & 1.0740 & 0.0 & 通过 \\
    S012 & 1 & 2.7524 & 0.0 & 通过 \\
    S013 & 1 & 1.8248 & 0.0 & 通过 \\
    S014 & 1 & 2.1404 & 0.0 & 通过 \\
    S015 & 1 & 2.3075 & 0.0 & 通过 \\
    \bottomrule
  \end{tabular*}
\end{table}

从算法规模看，单个服务区枚举的可行模式数最多为 201 个，记忆化动态规划实际访问的状态数最多为 165 个，因此问题一可以在很短时间内精确求解。程序记录的本次问题一总运行时间约为 $0.679\,\mathrm{s}$。该结果也说明，利用“同属性货箱归类 + 状态压缩”的方式能够显著减少离散组合规模。

\subsection{问题一求解结果总结与讨论}

针对单服务区直接往返运输能力与不可拆货箱组批问题，本文在不考虑实体无人机数量及共享电池调度的条件下，依据题设地形净空、飞行时间与运输能耗计算规则，建立最大安全载荷模型，并综合货箱质量、装载体积和返航安全能量约束，采用动态规划求解各服务区的货箱组批方案。针对问题一的各项要求，主要结果如下。

\noindent\textbf{（1）三种机型在不同服务区的最大安全载荷。}

在默认 $20\%$ 返航安全余量下，A 型无人机在全部 15 个服务区的最大安全载荷均为 $25\,\mathrm{kg}$；B 型在 S008 的最大安全载荷为 $28.801\,\mathrm{kg}$，其余服务区均为 $30\,\mathrm{kg}$；C 型在 S002、S003、S004、S008 和 S012 的最大安全载荷分别为 $68.860$、$68.206$、$63.697$、$58.904$ 和 $68.117\,\mathrm{kg}$，其余服务区均为 $80\,\mathrm{kg}$。

上述结果表明，山区单点往返任务的最大安全载荷不仅取决于机型额定载质量，还受到水平航程、地形引起的爬升需求及返航安全余量的共同约束。在当前参数下，A 型均可达到额定载荷上限，而 B、C 两型在部分服务区受到能量约束。其中，S008 同时使 B、C 两型的最大安全载荷低于额定上限，是本算例中具有代表性的受限服务区。因此，不能直接以额定载质量替代各服务区的实际安全载荷进行货箱组批。

\noindent\textbf{（2）满足货箱与运输约束的组批方案。}

在默认条件下，80 个货箱被划分为 18 个直接往返架次，每个货箱恰好安排一次，各架次均满足载质量、装载体积及返航安全能量约束。其中，B 型、C 型各承担 9 个架次，A 型未被采用。具体而言，S001 由两个 C 型架次完成，载荷分别为 $76\,\mathrm{kg}$ 和 $78\,\mathrm{kg}$；S002、S003 均由一个 B 型架次和一个 C 型架次完成，载荷分别为 $25\,\mathrm{kg}$ 和 $56\,\mathrm{kg}$；S004 至 S008 各由一个 C 型架次完成；S009 至 S015 各由一个 B 型架次完成。各架次的货箱组成、载荷、体积、能耗及返航剩余电量见前文组批结果表。

A 型未被选中并不意味着其无法完成安全往返，而是质量、体积约束与优化目标共同作用的结果。例如，由饮用水、医疗物资和应急食品构成的典型组合总质量为 $25\,\mathrm{kg}$，但总体积为 $0.067\,\mathrm{m}^{3}$，超过 A 型 $0.060\,\mathrm{m}^{3}$ 的装载体积，却未超过 B 型 $0.073\,\mathrm{m}^{3}$ 的装载体积。因此，即使货箱组合满足质量上限，也未必能够由相应机型单架次完成运输，组批时必须同时检查质量与体积约束。

\noindent\textbf{（3）架次数、运输能耗与累计作业时间的优化结果及权衡。}

本文采用“首先最小化架次数，其次最小化总运输能耗，最后最小化累计作业时间”的字典序优先关系。所得方案共需 18 个架次，总运输能耗为 $59.1313\,\mathrm{kWh}$，累计作业时间为 $9.104\,\mathrm{h}$。这里的累计作业时间是全部架次作业时长之和，并非考虑多架无人机并行执行后的任务完工时间。

对照实验表明，将能耗改为第一优先目标后，总运输能耗降至 $59.0339\,\mathrm{kWh}$，降幅约为 $0.165\%$；但架次数增加至 19，增幅约为 $5.56\%$，累计作业时间增加至 $9.567\,\mathrm{h}$，增幅约为 $5.08\%$。以累计作业时间为第一目标时，所得方案与本文采用的字典序方案一致。因此，在本算例中，架次数优先能够以较小的能耗代价减少运输批次数和累计作业时间。该结果说明本文方案是在明确优先关系下取得的折中解，而不是三个指标分别达到各自独立最优的方案。

\noindent\textbf{（4）返航安全余量变化对最大安全载荷与组批结果的影响。}

随着返航安全余量提高，可用于运输的能量预算减少。对于固定的机型与服务区，最大安全载荷随余量增加不会上升，而是可能保持额定上限、逐步下降，直至该组合无法满足空载安全往返要求。以 S008 为例，安全余量由 $20\%$ 提高至 $30\%$ 时，A、B、C 三型的最大安全载荷分别由 $25.000$、$28.801$、$58.904\,\mathrm{kg}$ 降至 $14.213$、$21.001$、$36.814\,\mathrm{kg}$；余量达到 $35\%$ 时，A 型已无法满足空载安全往返要求，C 型的最大安全载荷降至 $15.683\,\mathrm{kg}$。

在组批层面，已测试的 $10\%$、$15\%$ 和 $20\%$ 余量下，最优方案均保持 18 个架次，总运输能耗均为 $59.1313\,\mathrm{kWh}$。余量提高至 $25\%$、$30\%$ 和 $35\%$ 时，架次数分别增加至 19、20 和 25，总运输能耗分别增加至 $61.073$、$67.228$ 和 $75.608\,\mathrm{kWh}$，累计作业时间分别增加至 $9.602$、$10.163$ 和 $12.607\,\mathrm{h}$。当余量为 $40\%$ 时，在现有机型与单点往返规则下，已不存在能够完成全部货箱交付的可行组批方案。

综上，问题一给出了三种机型对 15 个服务区的最大安全载荷、满足约束的完整货箱组批方案、不同优化优先级下的指标权衡，以及返航安全余量对载荷能力和组批结果的影响。
\FloatBarrier

\subsection{算法示例}

数学建模求解算法示例：
\begin{center}
\begin{minipage}{0.8\textwidth}
\begin{algorithm}[H]%[!htp]
\caption{算法的名字} %算法的名字
{\bf 输入:} %算法的输入， \hspace*{0.02in}用来控制位置，同时利用 \\ 进行换行
input parameters A, B, C\\
{\bf 输出:} %算法的结果输出
output result
\begin{algorithmic}[1]
\State some description 算法介绍 % \State 后写一般语句
\For{condition} % For 语句，需要和EndFor对应
  \State ...
  \If{condition} % If 语句，需要和EndIf对应
    \State ...
  \Else
    \State ...
  \EndIf
\EndFor
\While{condition} % While语句，需要和EndWhile对应
  \State ...
\EndWhile
\State \Return result
\end{algorithmic}
\end{algorithm}
\end{minipage}
\end{center}
\vspace{2ex}

\subsection{表格}

使用 tabular 环境, 如下表格: 表~\ref{tab:foo}. 通过 \verb|autoref| 引用表格: \autoref{tab:foo}.

这里可使用命令 \verb|\tabcolsep| 和 \verb|\arraystretch| 分别控制列间距和行间距.

\begin{table}[htp!]
\centering
\setlength{\tabcolsep}{12pt}  % 6pt standard
\renewcommand{\arraystretch}{1.2}
\caption{学术活动安排样例}
\label{tab:foo}
\begin{tabular}{|c|c|c|c|}
\hline
\textbf{日期}  & \textbf{地点} & \textbf{活动名称} & \textbf{备注} \\ \hline
2024年8月1日      & 上海       & 学术研讨会      & 主题：人工智能 \\ \hline
2024年8月15日    & 北京       & 学术交流会      & 重点：数学建模 \\ \hline
2024年9月1日      & 深圳       & 研究研讨会      & 主题：数据科学 \\ \hline
2024年10月15日  & 广州       & 创新论坛         & 重点：科技创新 \\ \hline
\end{tabular}
\end{table}


\clearpage

三线表
\begin{table}[htp!]
\newcolumntype{L}{X}
\newcolumntype{C}{>{\centering \arraybackslash}X}
\newcolumntype{R}{>{\raggedright \arraybackslash}X}
\centering
\caption{某校学生升高体重样本}
\label{tab:heightweight}
\begin{tabularx}{0.9\textwidth}{CCCC}
\toprule
序号 & 年龄 & 身高 & 体重 \\
\midrule
001 & 15 & 156 & 42 \\
002 & 16 & 158 & 45 \\
003 & 14 & 162 & 48 \\
004 & 15 & 163 & 50 \\
\cmidrule{2-4}
平均 & 15 & 159.75 & 46.25 \\
\bottomrule
\end{tabularx}
\end{table}

% 某行业产量与生产费用的数据
% \begin{table}[htp!]
% \centering
% \caption{某行业产量与生产费用的数据}%\label{}
% \newcolumntype{Y}{>{\centering\arraybackslash}X}
% \newcolumntype{Z}{!{\vline}@{\color{white}\vrule width \doublerulesep}!{\vrule}}%自定义列格式(双线)
% \begin{tabularx}{0.94\textwidth}{c|c|YZc|c|Y}
% \Xhline{0.9pt}
% 企业编号&  产量(台)&生产费用(万元)&企业编号&产量(台)&生产费用(万元)\\\Xcline{1-3}{0.6pt}\Xcline{4-6}{0.6pt}
% 1 & 40 & 130 & 7&  84&  165\\
% 2 & 42 & 150 & 8&  100&  170\\
% 3 & 50 & 155 & 9&  116&  167\\
% 4 & 55 & 140 & 10&  125&  180\\
% 5 & 65 & 150 & 11&  130&  175\\
% 6 & 78 & 154 & 12&  140&  185\\\Xhline{0.72pt}
% \end{tabularx}
% \end{table}

研究生数学建模2019年F题结果示例
\begin{table}[htp!]
\centering
\caption{问题1结果1 (左) 与 问题2结果 (右)}
\begin{minipage}[t]{0.48\linewidth}
\renewcommand\arraystretch{1.2} %定义表格高度
\newcolumntype{Y}{>{\centering\arraybackslash}X}
\begin{tabularx}{0.9\textwidth}{|Y|Y|Y|}
  \hline
  数据集1  &  数据集1 & 数据集2  \\
 \hline
  A问题1   & A问题1 &  A问题1     \\
  \hline
  503   &  503     & 163     \\
  294   &  200     & 114      \\
  91    &  80      & 8     \\
  607   &  237     & 309      \\
  540   &  170     & 305    \\
  250   &  278     & 123    \\
  340   &  369     & 45      \\
  277   &  214     & 160    \\
  B     &  397     & 92    \\
        &  B       & 93    \\
        &          & 61        \\
        &          & 292       \\
        &          & B         \\
104861  & 103518   & 109342 \\
\hline
\end{tabularx}
\end{minipage}
\begin{minipage}[t]{0.48\linewidth}
\renewcommand\arraystretch{1.2} %定义表格高度
\newcolumntype{Y}{>{\centering\arraybackslash}X}
\begin{tabularx}{0.9\textwidth}{|Y|Y|Y|}
  \hline
  数据集1  &  数据集1 & 数据集2  \\
  \hline
  A问题2  &  A问题2  & A问题2    \\
  \hline
   503    & 503     & 163  \\
  294    & 200     & 114   \\
   91     & 80     & 8   \\
   607    & 237     & 309   \\
   540    & 170    & 305  \\
   250    & 278    & 123  \\
   340    & 369     & 45   \\
   277    & 214    & 160  \\
  B      & 397   & 92 \\
        &  B     &  93      \\
         &         &  61      \\
        &         &   292     \\
        &         &   B     \\
 104917  &  103563  &109427 \\
\hline
\end{tabularx}
\end{minipage}
\end{table}

\clearpage
\begin{table}[htp!]
%\small
\centering
\renewcommand\arraystretch{1.2} %定义表格高度
\newcolumntype{Y}{>{\centering\arraybackslash}X}
\caption{问题3结果}
\begin{tabularx}{0.9\textwidth}{|Y|Y|Y|Y|Y|Y|}
  \hline %定义表格宽度
  数据集1  &  数据集1  &  \multicolumn{2}{c|}{数据集2 (无问题点)}  & \multicolumn{2}{c|}{数据集2 (有问题点)} \\
  \hline
  A问题3   & A问题3   &  \multicolumn{2}{c|}{A问题3}    &  \multicolumn{2}{c|}{A问题3}      \\
  \hline
  503      & 503     & 169      & 73      & 169     & 73   \\
  \hline
  69       & 69      & 322      & 249     & 322     & 249   \\
  \hline
  506      & 506      & 270     & 274     & 270     & 274   \\
  \hline
  371      & 371      & 89      & 12      & 89      & 12   \\
  \hline
  183      & 183      & 236     & 216     & 236     & 216  \\
  \hline
  194      & 194      & 132     & 16      & 132     & 16   \\
  \hline
  450      & 450      & 53      & 282     & 53      & 282   \\
  \hline
  286      & 113      & 112     & 84      & 112     & 141  \\
  \hline
  485      & 485      &  268    & 287     & 268     & 291 \\
  \hline
 \red{B (9D)}~~     & 248      & 250     & 99      & 250     &161 \\
 \hline
   & \red{B (10D)}   & 243     & \red{B (21D)}   & 243    & \red{B (21D)} \\
   \hline
    &          &         &         &         &       \\
    \hline
  104861m  & 103518m  &         & 168924m  &         &161650m \\
\hline
\end{tabularx}
\end{table}


\clearpage

\section{问题二的建模与求解}

\subsection{问题二分析}
问题二在问题一的基础上，由单点往返扩展为\textbf{异构无人机多点、多架次联合调度}。其主要难点在于货箱组批、服务区访问顺序、无人机选择、架次开始时间以及共享电池周转相互耦合，若直接对所有决策变量同时搜索，将产生较大的组合规模。

因此，本问可采用\textbf{“可行架次生成 + 资源时序调度”}的分层求解思路。首先依据问题一得到的载荷与能耗计算方法，生成满足重量、体积、飞行能量和返航安全要求的候选运输架次，并确定每个架次可访问的服务区及其访问顺序；随后将这些候选架次分配给具体无人机和共享电池，并确定各架次开始时刻。调度过程中重点处理医疗物资、首批保障物资等具有较强时效要求的货箱，同时保证无人机任务占用、电池充电及再次使用之间不存在时间冲突。最终以\textbf{配送及时性、任务完成时间、总能耗和架次数}作为主要评价指标，形成兼顾救援效率与资源利用率的运输调度方案。


\section{问题三的建模与求解}
\subsection{问题三分析}
问题三的关键是将问题二中的运输调度进一步与\textbf{连续通信保障}相结合。山区环境中，运输无人机与固定网关之间的通信状态会随位置、高度和地形遮挡不断变化，因此不能仅判断服务区处是否具有通信条件，而应对无人机的整个飞行轨迹进行通信可用性检测。

首先根据运输无人机在爬升、巡航、下降和投送阶段的空间轨迹，结合DEM地形进行视距遮挡判断，并利用链路预算计算无人机与固定网关之间的通信状态。对于无法直连的轨迹区段，引入中继无人机提供通信保障。由于中继悬停位置属于连续空间变量，直接联合优化难度较大，因此可在满足地形和高度限制的区域内构造\textbf{候选中继位置集合}，预先计算各候选位置对不同运输轨迹和时间区段的通信覆盖关系，将连续空间问题转化为有限候选点上的时空覆盖问题。

在此基础上，将运输架次、中继位置、中继服务时段和能源组件进行联合调度，保证每一时刻运输无人机至少能够通过固定网关直连或中继链路保持通信。故本问的核心突破口在于\textbf{先建立轨迹级通信可行性矩阵，再进行运输--中继协同调度}，从而在保证通信连续性的同时，兼顾任务完成时间、运输与中继能耗以及设备使用规模。

\section{问题四的建模与求解}
\subsection{问题四分析}
问题四是在问题三既定联合调度方案的基础上研究\textbf{任务分区与资源独立配置}问题。由于题目要求保持原有货箱组批、访问顺序、运输任务和通信保障关系不变，因此本问并不重新规划运输路线，而是需要识别哪些服务区之间因共享同一运输架次或通信任务而具有不可分割的关联关系。

为此，可首先构建\textbf{服务区任务关联图}：将各服务区视为节点，若多个服务区出现在同一运输架次中，则在其之间建立关联，并将相互关联的服务区合并为不可拆分的任务单元。在此基础上，分别对这些任务单元进行两分区和三分区，使各组在任务量和资源需求上尽可能均衡。

完成分区后，需要根据每个任务组中的运输与中继任务时间区间，计算运输无人机、共享电池、中继无人机和能源组件的\textbf{最大并发占用量}，而不能简单累加任务总数，以此确定各分区独立运行时真正需要配置的最小资源数量。最终从资源配置规模、资源冗余、工作量均衡程度及现有库存缺口等方面比较两种分区方案，并识别导致资源不足的关键任务和时间段。

\section{模型评价}
\subsection{模型的优缺点}
\subsection{模型的改进方向}


%参考文献   手工录入
%\begin{thebibliography}{9}%宽度9
% \bibitem{bib:one} ....
% \bibitem{bib:two} ....
%\end{thebibliography}

%采用bibtex方案
\cite{mittelbach_latex_2004,wright_latex3_2009,beeton_unicode_2008,vieth_experiences_2009}

\bibliographystyle{gmcm}
\bibliography{reference}


\clearpage
%附录
\begin{appendices}
%\setcounter{page}{1} %如果需要可以自行重置页码。
\section{MATLAB 源程序}
\renewcommand{\thesubsection}{A\thinspace.\thinspace\arabic{subsection}}
\subsection{第1问程序}
\vspace{-2ex}

\begin{Matlab}{code.m}
clear all
kk=2;
[mdd,ndd]=size(dd);
while ~isempty(V)
    [tmpd,j]=min(W(i,V));
    tmpj=V(j);
    for k=2:ndd
        [tmp1,jj]=min(dd(1,k)+W(dd(2,k),V));
        tmp2=V(jj);
        tt(k-1,:)=[tmp1,tmp2,jj];
    end
    tmp=[tmpd,tmpj,j;tt];
    [tmp3,tmp4]=min(tmp(:,1));
    if tmp3==tmpd,
        ss(1:2,kk)=[i;tmp(tmp4,2)];
    else
        tmp5=find(ss(:,tmp4)~=0);
        tmp6=length(tmp5);
        if dd(2,tmp4)==ss(tmp6,tmp4)
            ss(1:tmp6+1,kk)=[ss(tmp5,tmp4);tmp(tmp4,2)];
        else, ss(1:3,kk)=[i;dd(2,tmp4);tmp(tmp4,2)];
        end
    end
    dd=[dd,[tmp3;tmp(tmp4,2)]];
    V(tmp(tmp4,3))=[];
    [mdd,ndd]=size(dd);kk=kk+1;
end;
S=ss; D=dd(1,:);
\end{Matlab}
\vspace{2ex}

\clearpage
\section{Python 源程序}
\renewcommand{\thesubsection}{B\thinspace.\thinspace\arabic{subsection}}
\subsection{第2问程序}
\vspace{-2ex}
\begin{Python}{mip1.py}
# This example formulates and solves the following simple MIP model:
#  maximize
#        x +   y + 2 z
#  subject to
#        x + 2 y + 3 z <= 4% !Mode:: "TeX:UTF-8"
%!TEX program  = xelatex

\documentclass[bwprint,fontset=fandol]{gmcmthesis}

\hypersetup{hidelinks}

%========================================================
% 题目、摘要、关键词标题：隶书，小二
% 字体文件位置：Fonts/SIMLI.TTF
%========================================================

% 定义真正的隶书字体
\setCJKfamilyfont{mylishu}[
    Path = Fonts/
]{SIMLI.TTF}

% 将模板中的 \lishu 指向上传的真正隶书
\renewcommand{\lishu}{\CJKfamily{mylishu}}

%========================================================
% 将“关键词”三个字统一修改为：隶书、小二
%========================================================
\makeatletter

\renewcommand{\keywords}[1]{%
    \renewcommand{\gmcm@tokens@keywords}{#1}%
    \par
    \vskip1ex
    {\noindent\zihao{-2}\lishu
    \gmcm@cap@keywordsname：}%
    ~{\gmcm@tokens@keywords}%
}

\makeatother
\usepackage{graphicx}   % 必须，否则 \includegraphics 报错

% Support for the added equations, editable tables and figure placeholders.
\usepackage{amsmath,amssymb}
\usepackage{array,booktabs,tabularx,longtable}
\usepackage{float,placeins}

\usepackage[framemethod=TikZ]{mdframed}

%\usepackage{fontspec}  % for Consolas & Courier New
%\lstset{basicstyle=\small\fontspec{Consolas}}
%\lstset{basicstyle=\small\fontspec{Courier New}}
%\setmonofont{Consolas}
%\setcounter{tocdepth}{3}   % 调整目录深度
\usepackage{siunitx}
\usepackage{colortbl}
\definecolor{color1}{rgb}{0.78,0.88,0.99}
\definecolor{color2}{rgb}{0.36,0.62,0.84}
%\definecolor{color3}{rgb}{0.8235,0.8706,0.9373}
\definecolor{color3}{rgb}{0.88,0.92,0.96}
\definecolor{color4}{rgb}{0.96,0.97,0.98}%{0.9176,0.9373,0.9686}

% 算法
\usepackage[noend]{algpseudocode}
\usepackage{algorithmicx,algorithm}
\floatname{algorithm}{算法}
\renewcommand{\algorithmicrequire}{\textbf{输入:}}
\renewcommand{\algorithmicensure}{\textbf{输出:}}

%\numberwithin{equation}{section}
%\numberwithin{figure}{section}
%\numberwithin{table}{section}


\newcommand{\red}[1]{\textcolor{red}{#1}}
\newcommand{\blue}[1]{\textcolor{blue}{#1}}


%===================== 建模论文题目 ======================

\title{山区洪涝灾害下无人机运输与通信协同优化
}
\baominghao{\hspace{6em} 26106350114} %参赛队号
\schoolname{\hspace{5.8em} 西南大学} %学校名称
\membera{\hspace{6em} 谢文哲} %队员A
\memberb{\hspace{6em} 章林怡} %队员B
\memberc{\hspace{6em} 刘宇豪} %队员C

%=======================================================


\begin{document}

%生成标题
\maketitle

%填写摘要
\begin{abstract}

山区洪涝灾害往往伴随着\textbf{道路阻断、通信受损与能源受限}等多重约束，使传统地面运输难以及时完成应急物资配送。无人机具有机动性强、受地面道路影响小等优势，但其实际救援能力同时受到载荷、地形、能耗、配送时限、通信覆盖以及设备周转等因素制约。针对山区洪涝灾害场景下无人机应急运输与通信保障的协同优化问题，本文结合服务区空间位置、数字高程模型（DEM）、货箱需求、异构无人机参数、电池资源及无线通信链路参数，建立了由\textbf{单点运输能力分析、多机多架次调度、运输--通信联合调度和任务分区资源配置}构成的递进式优化模型，并从运输效率、任务及时性、能源消耗和资源占用等多个维度对方案进行评价。

针对问题一，为准确刻画不同机型在山区地形条件下的运输能力，首先利用\textbf{DEM栅格采样}提取调度中心与各服务区直线路径上的最高地形高程，并按照“爬升--巡航--下降”的飞行过程建立时间与能耗计算模型。在载荷影响等效航程且返航需保留安全能量余量的条件下，利用\textbf{二分搜索}求解各“机型--服务区”组合的最大安全载荷。在此基础上，将不可拆货箱的组批过程转化为同时考虑载质量、装载体积与安全能量约束的\textbf{多维装箱优化问题}，以架次数、总运输能耗和累计作业时间构造分层多目标优化模型，获得各服务区的货箱组批与机型选择方案。同时改变返航安全余量进行敏感性分析，刻画安全冗余增加对最大可运输载荷及架次数的影响规律。

针对问题二，在单架次可访问多个服务区的条件下，进一步考虑\textbf{异构无人机、共享电池、配送时限与多架次周转}之间的耦合关系。本文构建带容量、能量和时间窗约束的异构多旅行车辆路径模型，将货箱组批、服务区访问顺序、运输机型、实体无人机、电池资源及架次开始时刻进行联合决策。为降低大规模组合搜索的复杂度，采用\textbf{“可行架次生成--资源调度优化”两阶段求解框架}：首先生成满足载荷、体积和能量限制的候选运输架次，再结合医疗物资期望送达时间和首批保障截止时间，通过自适应邻域搜索与整数优化对路线及资源时序进行协调。模型能够输出逐架次运输路线、逐箱送达时刻以及无人机和电池占用情况，并通过任务时序和资源冲突检验保证方案的实际可执行性。

针对问题三，在运输调度基础上进一步引入山区地形遮挡和无线传播约束。首先依据运输无人机的连续飞行轨迹，结合DEM完成\textbf{视距（LOS）遮挡判定}，并根据双向链路预算计算运输无人机与固定网关之间的实时通信可用性。对于直连链路失效区段，在DEM覆盖区域内生成满足悬停高度约束的候选中继位置，并预计算“运输无人机--中继无人机--固定网关”两段链路的可服务关系，从而将连续空间中继部署问题转化为有限候选位置上的\textbf{时空覆盖问题}。进一步建立运输无人机与中继无人机联合调度模型，同时优化运输架次、中继悬停位置、服务时段及能源组件配置，在保证运输无人机\textbf{全飞行过程连续通信}的前提下，综合降低物资延误、联合任务完成时间、运输与中继总能耗以及设备使用强度，并通过时间离散精度的敏感性分析验证通信判定结果的稳定性。

针对问题四，为研究灾后分区独立救援条件下的资源配置问题，根据问题三中已经确定的运输架次和通信保障关系构建\textbf{服务区任务耦合图}。若多个服务区位于同一运输架次或存在不可分割的联合保障关系，则将其合并为不可拆分的任务单元，并分别建立两分区和三分区优化模型。在保持原运输与中继任务安排不变的条件下，以\textbf{资源配置规模、资源冗余、组间工作量均衡和现有库存缺口}为主要评价指标进行分区优化；同时基于任务占用时间区间计算各分区运输无人机、共享电池、中继无人机及能源组件的最大并发需求，从而得到各分区独立执行所需的最小资源配置，并识别潜在资源瓶颈及其形成原因。

综上，本文将\textbf{地形约束、运输调度、能源周转与连续通信保障}统一纳入山区应急无人机救援决策框架，在保证货箱不可拆分、任务时限、返航安全以及通信连续性的基础上，实现了由单架次能力分析到多机协同、再到运输--通信联合优化与灾区资源分区配置的逐层扩展。所建立模型具有明确的物理含义和较好的可解释性，可为山区洪涝灾害条件下无人机应急运输与通信资源协同配置提供定量化决策依据。

\keywords{山区洪涝灾害\quad  无人机应急运输\quad 异构多架次调度\quad  通信中继  \quad  协同优化}
\end{abstract}

\pagestyle{plain}

%目录 不推荐加
\maketoc

\clearpage

\section{问题重述}

\subsection{引言}   % 问题的背景
山区洪涝灾害具有突发性强、地形复杂和基础设施易受损等特点。持续强降雨可能引发山洪、滑坡及道路塌方，并进一步造成道路、供电和通信设施同时中断，使部分受灾区域在短时间内难以通过传统地面交通获得物资补给。题目所描述的广西横州市镇龙乡洪涝灾害中，就曾出现断路、断电、断网以及大量群众受困的情况，食品、饮用水、药品及抢修设备需要依靠徒步或空中方式运入灾区。 因此，利用受地面道路条件影响较小的无人机开展应急物资运输，对提高山区灾害救援效率具有重要意义。

然而，山区无人机救援并非简单的物资配送问题。一方面，复杂地形会影响无人机的飞行高度、航程及能源消耗，不同货箱的重量、体积和配送时限也使货箱组批、路径规划、机型选择和电池周转相互制约；另一方面，山体遮挡可能导致运输无人机与固定网关之间通信中断，需要利用中继无人机提供持续通信保障。由此，运输任务、通信保障和有限装备资源之间形成了明显的耦合关系。题目也明确要求同时考虑载荷、地形净空、能源储备、任务时限及通信链路等多种因素。

基于此，本文以题目给出的调度中心、15个服务区、80个不可拆货箱以及运输无人机、中继无人机、共享电池和DEM等数据为基础，围绕单点运输能力、多机多架次调度、运输与通信协同以及任务分区资源配置四个层次展开建模，在保证运输安全、配送时效和通信连续性的前提下，对有限无人机及能源资源进行合理配置，为山区洪涝灾害下的无人机应急救援提供定量化决策依据。
\subsection{问题的提出}

%\subsubsection{问题的提出内容一}

\noindent 根据题目给定的调度中心、服务区、货箱需求、DEM地形数据、运输与中继无人机参数、电池资源及通信链路参数，需要依次解决以下四个相互递进的问题。

\textbf{问题一：} 单点往返运输能力与货箱组批。针对不同机型在各服务区执行
$O_{01}\rightarrow S_i\rightarrow O_{01}$
往返任务的场景，结合DEM地形、载荷和能耗约束，确定各机型的\textbf{最大安全载荷}。在货箱不可拆分的条件下，优化\textbf{货箱组批方案}，并分析返航安全余量变化对运输能力和组批结果的影响。

\textbf{问题二：} 异构无人机多点多架次运输调度。在允许单架次访问多个服务区的条件下，综合考虑载荷、配送时限、无人机数量及共享电池周转等约束，联合优化\textbf{货箱组批、访问顺序、无人机分配和架次时序}，以提高配送及时性并降低任务完成时间、能耗和架次数。

\textbf{问题三：} 通信约束下的运输与中继联合调度。在问题二的基础上引入山区地形遮挡和通信链路约束，当运输无人机与固定网关 $G_{01}$ 无法直连时，通过中继无人机提供通信保障。联合优化\textbf{运输任务与中继无人机的悬停位置、服务时段及能源配置}，保证运输全过程通信连续，并兼顾时效与能耗。

\textbf{问题四：} 救援任务分区与资源配置。在保持问题三运输与通信任务关系不变的条件下，将 $15$ 个服务区分别划分为 $2$ 个和 $3$ 个任务组。比较不同分区方案下的\textbf{资源需求、资源冗余、工作量均衡及库存缺口}，确定更合理的分区与资源配置方案。

\clearpage

\section{模型假设与符号说明}
\subsection{模型假设}
\subsection{符号说明}
\begin{table}[htp!]
\centering
\renewcommand\arraystretch{1.2} %定义表格高度
\newcolumntype{L}{>{\quad}X}
\newcolumntype{P}[1]{>{\centering\arraybackslash}p{#1}}
\newcolumntype{C}{>{\centering \arraybackslash}X}
\newcolumntype{R}{>{\raggedright \arraybackslash}X}
\begin{tabularx}{0.9\textwidth}{|P{1.8cm}|L|}
  \hline %\toprule
  符号    &   \quad 意义 \\
  \hline %\midrule
  $ O $  &  临时调度中心    \\
  \hline
  $ S_i $  &  第 $ i $ 个服务区，$ i=1,2,\ldots,15 $    \\
  \hline
  $ g $  &  运输无人机机型，$ g \in \{A,B,C\} $    \\
  \hline
  $ m_g $  &  机型 $ g $ 的无人机空机质量/kg    \\
  \hline
  $ Q_g $  &  机型 $ g $ 的额定最大载货质量/kg    \\
  \hline
  $ V_g $  &  机型 $ g $ 的最大装载体积/m$^3$    \\
  \hline
  $ E_g $  &  机型 $ g $ 单组电池可用能量/kWh    \\
  \hline
  $ p_g $  &  机型 $ g $ 的返航安全余量比例    \\
  \hline
  $ v_g $  &  机型 $ g $ 的水平巡航速度/(m·s$^{-1}$)    \\
  \hline
  $ v_g^\uparrow, v_g^\downarrow $  &  机型 $ g $ 的爬升速度和下降速度/(m·s$^{-1}$)    \\
  \hline
  $ \eta_g $  &  爬升附加能耗效率    \\
  \hline
  $ R_g^0, R_g^F $  &  机型 $ g $ 的空载、满载标准航程/m    \\
  \hline
  $ d_{ij} $  &  节点 $ i $ 到节点 $ j $ 的水平测地距离/m    \\
  \hline
  $ H_{ij} $  &  航线 $ i \to j $ 的计划巡航海拔/m    \\
  \hline
  $ h_{ij}^\uparrow, h_{ij}^\downarrow $  &  航线 $ i \to j $ 的爬升高度、下降高度/m    \\
  \hline
  $ q $  &  当前有效载荷/kg    \\
  \hline
  $ R_g(q) $  &  载荷 $ q $ 下的等效航程/m    \\
  \hline
  $ E_{ij}(g|q) $  &  机型 $ g $ 在载荷 $ q $ 下执行航线 $ i \to j $ 的运输能耗/kWh    \\
  \hline
  $ t_{ij}(g) $  &  机型 $ g $ 执行航线 $ i \to j $ 的飞行时间/s    \\
  \hline
  $ B_i $  &  服务区 $ S_i $ 的货箱集合    \\
  \hline
  $ m_b, v_b $  &  货箱 $ b $ 的质量/kg 与体积/m$^3$    \\
  \hline
  $ q_g^{\mathrm{safe}} $  &  机型 $ g $ 服务 $ S_j $ 时的最大安全载荷/kg    \\
  \hline
  $ p $  &  某一可行班驶模式    \\
  \hline
  $ c_p $  &  班驶模式 $ p $ 对各类货箱的数量向量    \\
  \hline
  $ N $  &  全部服务区总运输频次    \\
  \hline
  $ E $  &  全部班次总运输能耗/kWh    \\
  \hline
  $ T $  &  全部班次累计作业时间/s    \\
  \hline
\end{tabularx}
\end{table}


%\begin{tabularx}{\textwidth-18pt}{XXX}
%\hline
%Input & Output& Action return \\
%\hline
%DNF &  simulation & jsp\\
%\hline
%\end{tabularx}


\section{问题一的建模与求解}

\subsection{问题一分析}
问题一要求在不考虑实体无人机和共享电池调度的条件下，研究“临时调度中心 $O01$---单个服务区 $S_i$---$O01$”的直接往返运输能力与货箱组批问题。由于每一架次只能服务一个服务区，15 个服务区之间不存在跨区组批，因此该问题可以先对每个“机型—服务区”组合建立独立的飞行时间与能耗模型，再在各服务区内部求解不可拆货箱的最优组批，最后汇总得到全局结果。
与普通装箱问题相比，本题的关键难点在于：第一，载荷上限不仅由机型额定载质量决定，还受 DEM 地形引起的爬升能耗、航程衰减和返航安全余量共同约束；第二，货箱不可拆分，同时受到质量和体积双重限制；第三，题目要求综合考虑往返架次数、运输能耗与累计作业时间，三者之间存在明显的层级权衡。因此，本文采用“物理可行性计算—可行组批模式枚举—状态压缩动态规划”的思路进行精确求解，整体流程如图~\ref{fig:problem1_flow}所示。
\begin{figure}[htbp!]
  \centering
  \includegraphics[width=0.85\textwidth]{figures/Q1_FIg1.png}
  \caption{问题一建模与求解总框架}
  \label{fig:problem1_flow}
\end{figure}
\subsection{基本约定}
（1）每一运输架次均采用 $O01 \to S_i \to O01$ 的直接往返形式，只服务一个服务区；同一服务区允许使用多个架次完成配送，但不得跨服务区组批。

（2）任意两个任务节点之间的水平航迹取两节点连线；巡航海拔由该连线经过的 $30\,\mathrm{m}$ DEM 像元最高地形高程加 $50\,\mathrm{m}$ 安全净空确定。

（3）$O01$ 的作业高度取地面海拔，服务区作业高度取地面海拔以上 $30\,\mathrm{m}$；每个航段均按"爬升—水平巡航—下降"三个阶段计算飞行时间。

（4）下降阶段不单独计入附加势能消耗；运输总能耗由载荷相关水平航程能耗与爬升附加能耗构成。

（5）问题一不考虑实体无人机数量、共享电池数量及其充电周转，只考察单架次物理可行性与货箱组批；货箱不可拆分且每个货箱只能被安排一次。
% === Problem one: imported from the reference DOCX, sections 3.3--3.8 ===
% Data and conclusions are retained from the reference; no recomputation.
% Notation follows the existing symbol table: rho_g -> p_g,
% E_ij^g(q) -> E_ij(g|q), t_ij^g -> t_ij(g).
% The safe-payload symbol includes i to distinguish service areas.
\FloatBarrier

\subsection{航段时间与能耗模型}

\subsubsection{基于 DEM 的航段净空模型}

对于节点 $i$ 到节点 $j$ 的运输航段，沿节点连线逐个检查其穿越的 DEM 像元。设这些像元的最高地形高程为 $\max(h_k)$，则计划巡航海拔为

\begin{equation}
H_{ij} = \max_{k\in L_{ij}} h_k + 50.
\label{eq:q1-clearance}
\end{equation}

其中，$L_{ij}$ 表示航段 $i\to j$ 穿越的 DEM 像元集合。设节点 $i$、$j$ 的实际作业海拔分别为 $z_i^o$、$z_j^o$，则爬升高度和下降高度分别为

\begin{equation}
h_{ij}^{\uparrow} = H_{ij}-z_i^{o},\qquad
h_{ij}^{\downarrow} = H_{ij}-z_j^{o}.
\label{eq:q1-vertical-height}
\end{equation}

该处理方式能够将山区地形起伏直接映射到无人机的爬升需求和能量消耗中，避免仅按平面距离估计运输能力。

\subsubsection{载荷相关等效航程与运输能耗}

题目给出了各机型空载标准航程与满载标准航程。为描述载荷增加导致的航程衰减，采用题目配套计算规则中的载荷相关等效航程关系：

\begin{equation}
R_g(q)=R_g^0-\bigl(R_g^0-R_g^F\bigr)
\left(\frac{q}{Q_g}\right)^{3/2},\qquad 0\le q\le Q_g.
\label{eq:q1-range}
\end{equation}

于是，机型 $g$ 携带载荷 $q$ 执行航段 $i\to j$ 时，水平巡航能耗可按“可用电池能量$\times$航段距离/等效航程”折算；同时考虑爬升势能，航段总运输能耗为

\begin{equation}
E_{ij}(g\mid q)
= E_g\frac{d_{ij}}{R_g(q)}
+\frac{(m_g+q)g_0h_{ij}^{\uparrow}}{\eta_g\,3.6\times 10^6}.
\label{eq:q1-segment-energy}
\end{equation}

式中 $g_0=9.81\,\mathrm{m}/\mathrm{s}^{2}$；第二项除以 $3.6\times10^6$ 将焦耳换算为 kWh。由于题目规定下降能耗效率取 0，故下降阶段不再单独增加附加能耗。

\subsubsection{飞行时间与单架次作业时间}

机型 $g$ 在航段 $i\to j$ 上的飞行时间由爬升、巡航和下降三个阶段构成：

\begin{equation}
t_{ij}(g)
=\frac{h_{ij}^{\uparrow}}{v_g^{\uparrow}}
+\frac{d_{ij}}{v_g}
+\frac{h_{ij}^{\downarrow}}{v_g^{\downarrow}}.
\label{eq:q1-flight-time}
\end{equation}

对于问题一中的直接往返架次，若某批次含 $n$ 个货箱，则该架次累计作业时间还包括起飞准备、逐箱装载、服务区作业以及逐箱卸载。因此

\begin{equation}
\begin{aligned}
T_p={}&t_g^{\mathrm{prep}}+n\,t_g^{\mathrm{load}}
+t_{\mathrm{O01},i}(g)+t_{i,\mathrm{O01}}(g)\\
&+t_g^{\mathrm{service}}+n\,t_g^{\mathrm{unload}}.
\end{aligned}
\label{eq:q1-operation-time}
\end{equation}

\subsection{单点往返最大安全载荷模型}

问题一首先要求计算三种机型在各服务区执行直接往返任务时的最大安全载荷。起飞时无人机携带全部货物，到达服务区后完成投送，返航阶段按空载计算。因此，机型 $g$ 向服务区 $S_i$ 运送载荷 $q$ 时的一次往返总能耗为

\begin{equation}
E_{gi}^{\mathrm{rt}}(q)
=E_{\mathrm{O01},i}(g\mid q)+E_{i,\mathrm{O01}}(g\mid 0).
\label{eq:q1-roundtrip-energy}
\end{equation}

为保证返航安全，任务结束时必须至少保留 $p_g$ 比例的可用电量，即

\begin{equation}
E_{gi}^{\mathrm{rt}}(q)\le (1-p_g)E_g.
\label{eq:q1-energy-limit}
\end{equation}

因此，机型 $g$ 在服务区 $S_i$ 的最大安全载荷定义为

\begin{equation}
q_{gi}^{\mathrm{safe}}=\max\left\{q\;\middle|\;
0\le q\le Q_g,\quad E_{gi}^{\mathrm{rt}}(q)\le (1-p_g)E_g
\right\}.
\label{eq:q1-safe-payload}
\end{equation}

若额定载荷 $Q_g$ 已满足式~\eqref{eq:q1-energy-limit}，则 $q_{gi}^{\mathrm{safe}}=Q_g$；否则在 $[0,Q_g]$ 上利用单调求根得到使式~\eqref{eq:q1-energy-limit}恰好取等号的载荷。若 $q=0$ 时仍不能满足返航余量，则判定该“机型—服务区”组合在当前安全余量下不可达。
\begin{figure}[H]
  \centering
  \includegraphics[width=0.85\textwidth]{figures/payload_distance.pdf}

  \caption{典型困难服务区 S008 的安全载荷随返航余量变化}
  \label{fig:q1-s008-sensitivity}
\end{figure}

\subsection{货箱组批优化模型}

\subsubsection{可行组批约束}

对服务区 $S_i$ 的货箱集合 $B_i$，任意候选批次 $p\subseteq B_i$ 只能选择一种机型 $g_p$。设该批次的货箱总质量和总体积分别为 $M_p$ 与 $V_p$，则必须同时满足

\begin{equation}
M_p=\sum_{b\in p}m_b\le Q_{g_p},\qquad
V_p=\sum_{b\in p}v_b\le V_{g_p}.
\label{eq:q1-capacity}
\end{equation}

此外，该批次还必须满足返航安全能量约束

\begin{equation}
E_{\mathrm{O01},i}(g_p\mid M_p)
+E_{i,\mathrm{O01}}(g_p\mid 0)
\le (1-p_{g_p})E_{g_p}.
\label{eq:q1-batch-energy}
\end{equation}

式~\eqref{eq:q1-capacity}和式~\eqref{eq:q1-batch-energy}共同决定一个货箱组合是否可以由某一机型完成。需要特别指出，最大安全载荷只约束“质量”这一连续变量，实际组批仍可能首先受到装载体积限制，因此不能简单依据载重上限进行贪心装箱。

\subsubsection{多目标优化与优先关系}

题目要求在完成全部货箱交付的前提下，综合考虑往返架次数、总运输能耗和累计作业时间。考虑到应急运输中减少架次数能够直接降低任务组织复杂度和设备切换需求，本文采用字典序优化，将架次数作为第一优先级，在最少架次方案集合内继续最小化运输能耗，最后最小化累计作业时间：

\begin{equation}
\min_{\mathrm{lex}}\;(N,E,T).
\label{eq:q1-objective}
\end{equation}

其中 $N=\sum_p x_p$，$E=\sum_p E_p x_p$，$T=\sum_p T_p x_p$；$x_p\in\{0,1\}$ 表示可行组批模式 $p$ 是否被选用。每个货箱必须且只能出现一次，即

\begin{equation}
\sum_{p:\,b\in p}x_p=1,\qquad \forall b\in B_i.
\label{eq:q1-exact-cover}
\end{equation}

\subsection{精确动态规划求解方法}

由于问题一禁止跨服务区组批，因此 15 个服务区可以分别精确求解。为降低状态维度，将同一服务区内“物资类别、单箱质量、单箱体积”完全相同的货箱归为一类。设该服务区共有 $K$ 类货箱，第 $k$ 类货箱数量为 $n_k$，定义状态向量 $\boldsymbol{x}=(x_1,x_2,\ldots,x_K)$，其中 $x_k$ 表示尚未分配的第 $k$ 类货箱数量。

首先枚举所有满足质量、体积和返航余量约束的可行组批模式 $p$，并记录其货箱数量向量 $\boldsymbol{c}_p$、机型 $g_p$、能耗 $E_p$ 和作业时间 $T_p$。令 $F(\boldsymbol{x})$ 表示完成状态 $\boldsymbol{x}$ 中全部剩余货箱所需的最优代价向量，则有递推关系

\begin{equation}
F(\boldsymbol{x})=\min_{\mathrm{lex}}\left\{
(1,E_p,T_p)+F(\boldsymbol{x}-\boldsymbol{c}_p)
\;\middle|\;
\boldsymbol{0}<\boldsymbol{c}_p\le\boldsymbol{x},\quad
p\,\text{可行}
\right\},
\label{eq:q1-dp}
\end{equation}

\begin{equation}
F(\boldsymbol{0})=(0,0,0).
\label{eq:q1-dp-initial}
\end{equation}

采用记忆化搜索保存已计算状态，能够避免对同一剩余货箱状态重复求解。各服务区求得的最优组批结果彼此独立，其架次数、能耗和作业时间可以直接相加，从而得到问题一的全局字典序最优解。为进一步验证动态规划结果，本文又以相同的可行组批模式建立 0-1 整数规划进行交叉检查，15 个服务区的最少架次数和最小能耗均与动态规划结果一致。

\subsection{计算结果与分析}

\subsubsection{三种机型的最大安全载荷}

在题目默认的 20\% 返航安全余量下，三种机型对 15 个服务区的最大安全载荷如图~\ref{fig:q1-safe-payload} 和表~\ref{tab:q1-safe-payload} 所示。A 型无人机在全部服务区均可达到其 $25\,\mathrm{kg}$ 的额定载荷；B 型除 S008 外均可达到 $30\,\mathrm{kg}$，S008 的最大安全载荷降至 $28.801\,\mathrm{kg}$；C 型受地形和航程影响最为明显，在 S002、S003、S004、S008 和 S012 的最大安全载荷分别为 68.860、68.206、63.697、58.904 和 $68.117\,\mathrm{kg}$，其余服务区仍可达到 $80\,\mathrm{kg}$ 的额定上限。

% Figure placeholder only; no external image is required.
\begin{figure}[H]
  \centering
\includegraphics[width=0.85\textwidth]{figures/3.2.pdf}

  \caption{最大安全载荷热力图}
  \label{fig:q1-safe-payload}
\end{figure}

\begin{table}[H]
  \centering
  \small
  \setlength{\tabcolsep}{5pt}
  \renewcommand{\arraystretch}{1.15}
  \caption{20\%返航安全余量下各服务区最大安全载荷（kg）}
  \label{tab:q1-safe-payload}
  \begin{tabular*}{0.75\linewidth}{@{\extracolsep{\fill}}crrr@{}}
    \toprule
    服务区 & A型 & B型 & C型 \\
    \midrule
    S001 & 25.000 & 30.000 & 80.000 \\
    S002 & 25.000 & 30.000 & 68.860 \\
    S003 & 25.000 & 30.000 & 68.206 \\
    S004 & 25.000 & 30.000 & 63.697 \\
    S005 & 25.000 & 30.000 & 80.000 \\
    S006 & 25.000 & 30.000 & 80.000 \\
    S007 & 25.000 & 30.000 & 80.000 \\
    S008 & 25.000 & 28.801 & 58.904 \\
    S009 & 25.000 & 30.000 & 80.000 \\
    S010 & 25.000 & 30.000 & 80.000 \\
    S011 & 25.000 & 30.000 & 80.000 \\
    S012 & 25.000 & 30.000 & 68.117 \\
    S013 & 25.000 & 30.000 & 80.000 \\
    S014 & 25.000 & 30.000 & 80.000 \\
    S015 & 25.000 & 30.000 & 80.000 \\
    \bottomrule
  \end{tabular*}
\end{table}

从结果可见，最大安全载荷并非仅由额定载质量决定。对 A 型而言，额定载荷较小，因此 20\% 余量下主要受到 $25\,\mathrm{kg}$ 载重上限约束；对重载 C 型而言，额定载荷虽然达到 $80\,\mathrm{kg}$，但在距离较远或需要较大爬升的服务区，载荷增加会同时降低等效航程并增加爬升势能消耗，因此安全载荷明显下降。S008 对 B、C 两型均形成约束，是后续敏感性分析中具有代表性的困难服务区。

\subsubsection{最优货箱组批方案}

在 20\% 返航安全余量和“架次数—能耗—累计作业时间”字典序目标下，80 个货箱共被划分为 18 个直接往返架次，总运输能耗为 $59.1313\,\mathrm{kWh}$，累计作业时间为 $9.104\,\mathrm{h}$，其中纯飞行时间为 $5.358\,\mathrm{h}$。最终使用 B 型 9 架次、C 型 9 架次，A 型未被采用。各服务区的组批结构如图~\ref{fig:q1-batching} 所示，完整架次结果见表~\ref{tab:q1-batching}。

% Figure placeholder only; no external image is required.
\begin{figure}[H]
  \centering
  
  \includegraphics[width=0.85\textwidth]{figures/flight_loading_red_blue.pdf}

  \caption{18 架次货箱组成图}
  \label{fig:q1-batching}
\end{figure}

\FloatBarrier
% The sortie table can span pages; column headers repeat automatically.
\begingroup
\small
\setlength{\tabcolsep}{4pt}
\renewcommand{\arraystretch}{1.2}
\setlength{\LTleft}{\fill}
\setlength{\LTright}{\fill}
\begin{longtable}{@{}lcclrrrrr@{}}
  \caption{问题一最优组批与运输结果}\label{tab:q1-batching}\\
  \toprule
    架次 & 服务区 & 机型 & 货箱组成 & \shortstack{载荷\\ /$\mathrm{kg}$} & \shortstack{体积\\ /$\mathrm{m}^{3}$} & \shortstack{能耗\\ /$\mathrm{kWh}$} & \shortstack{作业时间\\ /$\mathrm{min}$} & \shortstack{返航SOC\\ /$\%$} \\
  \midrule
  \endfirsthead
  \multicolumn{9}{c}{表~\thetable（续）}\\
  \toprule
    架次 & 服务区 & 机型 & 货箱组成 & \shortstack{载荷\\ /$\mathrm{kg}$} & \shortstack{体积\\ /$\mathrm{m}^{3}$} & \shortstack{能耗\\ /$\mathrm{kWh}$} & \shortstack{作业时间\\ /$\mathrm{min}$} & \shortstack{返航SOC\\ /$\%$} \\
  \midrule
  \endhead
  \midrule
  \multicolumn{9}{r}{\footnotesize 续下页}\\
  \endfoot
  \bottomrule
  \endlastfoot
    Q1-001 & S001 & C & \shortstack[l]{\strut $\mathrm{FOD}\times 2$\quad $\mathrm{HYG}\times 2$\\\strut $\mathrm{MED}\times 2$\quad $\mathrm{WAT}\times 3$} & 76 & 0.231 & 2.918 & 28.20 & 63.5 \\[3pt]
    Q1-002 & S001 & C & \shortstack[l]{\strut $\mathrm{FOD}\times 1$\quad $\mathrm{WAT}\times 5$} & 78 & 0.163 & 2.996 & 24.90 & 62.6 \\[3pt]
    Q1-003 & S002 & B & \shortstack[l]{\strut $\mathrm{FOD}\times 1$\quad $\mathrm{MED}\times 1$\\\strut $\mathrm{WAT}\times 1$} & 25 & 0.067 & 2.713 & 30.28 & 32.2 \\[3pt]
    Q1-004 & S002 & C & \shortstack[l]{\strut $\mathrm{FOD}\times 1$\quad $\mathrm{HYG}\times 1$\\\strut $\mathrm{WAT}\times 3$} & 56 & 0.144 & 5.721 & 34.03 & 28.5 \\[3pt]
    Q1-005 & S003 & B & \shortstack[l]{\strut $\mathrm{FOD}\times 1$\quad $\mathrm{MED}\times 1$\\\strut $\mathrm{WAT}\times 1$} & 25 & 0.067 & 2.780 & 34.68 & 30.5 \\[3pt]
    Q1-006 & S003 & C & \shortstack[l]{\strut $\mathrm{FOD}\times 1$\quad $\mathrm{HYG}\times 1$\\\strut $\mathrm{WAT}\times 3$} & 56 & 0.144 & 5.765 & 39.50 & 27.9 \\[3pt]
    Q1-007 & S004 & C & \shortstack[l]{\strut $\mathrm{FOD}\times 1$\quad $\mathrm{HYG}\times 1$\\\strut $\mathrm{MED}\times 1$\quad $\mathrm{WAT}\times 3$} & 59 & 0.156 & 6.153 & 38.35 & 23.1 \\[3pt]
    Q1-008 & S005 & C & \shortstack[l]{\strut $\mathrm{FOD}\times 1$\quad $\mathrm{HYG}\times 1$\\\strut $\mathrm{MED}\times 1$\quad $\mathrm{WAT}\times 3$} & 59 & 0.156 & 4.222 & 31.29 & 47.2 \\[3pt]
    Q1-009 & S006 & C & \shortstack[l]{\strut $\mathrm{FOD}\times 1$\quad $\mathrm{HYG}\times 1$\\\strut $\mathrm{MED}\times 1$\quad $\mathrm{WAT}\times 3$} & 59 & 0.156 & 2.598 & 26.03 & 67.5 \\[3pt]
    Q1-010 & S007 & C & \shortstack[l]{\strut $\mathrm{FOD}\times 1$\quad $\mathrm{HYG}\times 1$\\\strut $\mathrm{MED}\times 1$\quad $\mathrm{WAT}\times 2$} & 45 & 0.129 & 3.411 & 31.79 & 57.4 \\[3pt]
    Q1-011 & S008 & C & \shortstack[l]{\strut $\mathrm{FOD}\times 1$\quad $\mathrm{HYG}\times 1$\\\strut $\mathrm{MED}\times 1$\quad $\mathrm{WAT}\times 2$} & 45 & 0.129 & 5.836 & 36.66 & 27.0 \\[3pt]
    Q1-012 & S009 & B & \shortstack[l]{\strut $\mathrm{FOD}\times 1$\quad $\mathrm{MED}\times 1$\\\strut $\mathrm{WAT}\times 1$} & 25 & 0.067 & 2.096 & 25.97 & 47.6 \\[3pt]
    Q1-013 & S010 & B & \shortstack[l]{\strut $\mathrm{FOD}\times 1$\quad $\mathrm{MED}\times 1$\\\strut $\mathrm{WAT}\times 1$} & 25 & 0.067 & 1.823 & 25.92 & 54.4 \\[3pt]
    Q1-014 & S011 & B & \shortstack[l]{\strut $\mathrm{FOD}\times 1$\quad $\mathrm{MED}\times 1$\\\strut $\mathrm{WAT}\times 1$} & 25 & 0.067 & 1.074 & 19.81 & 73.2 \\[3pt]
    Q1-015 & S012 & B & \shortstack[l]{\strut $\mathrm{FOD}\times 1$\quad $\mathrm{MED}\times 1$\\\strut $\mathrm{WAT}\times 1$} & 25 & 0.067 & 2.752 & 32.05 & 31.2 \\[3pt]
    Q1-016 & S013 & B & \shortstack[l]{\strut $\mathrm{FOD}\times 1$\quad $\mathrm{MED}\times 1$\\\strut $\mathrm{WAT}\times 1$} & 25 & 0.067 & 1.825 & 25.17 & 54.4 \\[3pt]
    Q1-017 & S014 & B & \shortstack[l]{\strut $\mathrm{FOD}\times 1$\quad $\mathrm{MED}\times 1$\\\strut $\mathrm{WAT}\times 1$} & 25 & 0.067 & 2.140 & 31.16 & 46.5 \\[3pt]
    Q1-018 & S015 & B & \shortstack[l]{\strut $\mathrm{FOD}\times 1$\quad $\mathrm{MED}\times 1$\\\strut $\mathrm{WAT}\times 1$} & 25 & 0.067 & 2.307 & 30.47 & 42.3 \\[3pt]
\end{longtable}
\endgroup

注：同一服务区内同类别货箱的质量与体积参数一致，因此“货箱组成”按类别计数表示；FOD、HYG、MED、WAT 分别表示应急食品、生活卫生用品、医疗物资和饮用水。

结果显示，只有 S001、S002 和 S003 需要 2 个架次，其余 12 个服务区均可在 1 个架次内完成。其中 S001 总需求质量为 $154\,\mathrm{kg}$，必须由两架次 C 型承担；S002、S003 的总需求质量均为 $81\,\mathrm{kg}$，最优方案均采用“$25\,\mathrm{kg}$ 的 B 型架次 + $56\,\mathrm{kg}$ 的 C 型架次”。S004—S008 主要由 C 型单架次完成，S009—S015 则均由 B 型单架次完成。

值得注意的是，A 型无人机在本方案中没有被选中，这并不意味着其航程不可行，而是质量与体积约束共同作用的结果。以多个服务区中常见的“饮用水 $14\,\mathrm{kg}$ + 医疗物资 $3\,\mathrm{kg}$ + 应急食品 $8\,\mathrm{kg}$”组合为例，总质量恰为 $25\,\mathrm{kg}$，但总体积为 $0.067\,\mathrm{m}^{3}$，超过 A 型 $0.060\,\mathrm{m}^{3}$ 的装载体积，却可以装入 B 型 $0.073\,\mathrm{m}^{3}$ 的货舱。因此，若只按载重进行组批会错误高估 A 型的可用性。
\begin{figure}[H]
  \centering
  \fbox{%
    \begin{minipage}[c][4.2cm][c]{0.88\linewidth}
      \centering
      \small 图片占位框\par\medskip
      能耗优先方案相对本文方案的指标变化
    \end{minipage}%
  }
  \caption{载荷—距离—能耗响应曲面}
  \label{fig:q1-objective-comparison}
\end{figure}

\subsubsection{架次数、能耗与作业时间的权衡}

为检验字典序优先关系是否合理，进一步分别以能耗和累计作业时间作为第一目标进行对照求解，结果见表~\ref{tab:q1-objective-comparison}。

\begin{table}[H]
  \centering
  \small
  \setlength{\tabcolsep}{5pt}
  \renewcommand{\arraystretch}{1.15}
  \caption{不同优化优先级下的结果对比}
  \label{tab:q1-objective-comparison}
  \begin{tabular*}{\linewidth}{@{\extracolsep{\fill}}lrrrrr@{}}
    \toprule
    优化策略 & 架次数 & \shortstack{总能耗\\ /$\mathrm{kWh}$} & \shortstack{累计作业时间\\ /$\mathrm{h}$} & B型架次 & C型架次 \\
    \midrule
    架次数优先（本文） & 18 & 59.1313 & 9.104 & 9 & 9 \\
    能耗优先 & 19 & 59.0339 & 9.567 & 11 & 8 \\
    时间优先 & 18 & 59.1313 & 9.104 & 9 & 9 \\
    \bottomrule
  \end{tabular*}
\end{table}

% Figure placeholder only; no external image is required.
\begin{figure}[H]
  \centering
  
  \includegraphics[width=0.85\textwidth]{figures/运输与能耗变化率对比图.pdf}

  \caption{优化优先关系的代价比较}
  \label{fig:q1-objective-comparison}
\end{figure}

由表~\ref{tab:q1-objective-comparison} 可知，若将能耗置于第一优先级，总能耗可由 $59.1313\,\mathrm{kWh}$ 降至 $59.0339\,\mathrm{kWh}$，仅降低 0.165\%；但运输架次数由 18 增至 19，增加 5.56\%，累计作业时间也增加 5.08\%。而以累计作业时间为第一目标时，所得方案与本文字典序方案一致。可见，在本算例中，以减少架次数作为第一优先级只付出了极小的能耗代价，却能够降低任务批次数和总体作业时间，因而更符合灾后应急运输的组织需求。

\subsubsection{返航安全余量敏感性分析}

为分析安全裕度变化对运输能力和组批结果的影响，将返航安全余量从 10\% 逐步提高到 40\%，并在每个余量水平下重新计算安全载荷和最优组批。整体结果见表~\ref{tab:q1-reserve-sensitivity} 和图~\ref{fig:q1-reserve-sensitivity}。

\begin{table}[H]
  \centering
  \small
  \setlength{\tabcolsep}{5pt}
  \renewcommand{\arraystretch}{1.15}
  \caption{不同返航安全余量下的最优组批结果}
  \label{tab:q1-reserve-sensitivity}
  \begin{tabular*}{0.96\linewidth}{@{\extracolsep{\fill}}ccrrr@{}}
    \toprule
    返航余量 & 可行性 & 总架次数 & \shortstack{总能耗\\ /$\mathrm{kWh}$} & \shortstack{累计作业时间\\ /$\mathrm{h}$} \\
    \midrule
    10\% & 可行 & 18 & 59.131 & 9.104 \\
    15\% & 可行 & 18 & 59.131 & 9.104 \\
    20\% & 可行 & 18 & 59.131 & 9.104 \\
    25\% & 可行 & 19 & 61.073 & 9.602 \\
    30\% & 可行 & 20 & 67.228 & 10.163 \\
    35\% & 可行 & 25 & 75.608 & 12.607 \\
    40\% & 不可行 & - & - & - \\
    \bottomrule
  \end{tabular*}
\end{table}

当返航余量处于 10\%—20\% 时，虽然部分“机型—服务区”的连续安全载荷已经逐渐下降，但尚未跨越当前最优货箱组合的关键阈值，因此最优组批始终保持 18 架次，总能耗也不发生变化。当余量提高至 25\%、30\% 和 35\% 后，一些原本可以合并的货箱组合依次失去可行性，总架次数分别跃升至 19、20 和 25 架次，总能耗也由 $59.131\,\mathrm{kWh}$ 增长至 61.073、67.228 和 $75.608\,\mathrm{kWh}$。余量达到 40\% 时，全部货箱已经无法在现有三类机型下完成配送，整体问题变为不可行。

这一结果说明，由于货箱不可拆分，返航余量对组批方案的影响具有明显的“阈值效应”：安全载荷本身随余量连续下降，但架次数并不会连续变化，而是在安全载荷跌破某一离散货箱组合质量时发生跳变。

% Figure placeholder only; no external image is required.
\begin{figure}[H]
  \centering
  \includegraphics[width=0.85\textwidth]{figures/返航余量与组批结构、能耗、时间.png}
  \caption{返航余量与组批结构、能耗、时间}
  \label{fig:q1-s008-sensitivity}
\end{figure}

以 S008 为例，20\% 返航余量下 A、B、C 三型的最大安全载荷分别为 25.000、28.801 和 $58.904\,\mathrm{kg}$；余量提高到 30\% 后分别降至 14.213、21.001 和 $36.814\,\mathrm{kg}$。继续提高到 35\% 时，A 型即使空载也无法满足返航余量，而 C 型最大安全载荷仅剩 $15.683\,\mathrm{kg}$；到 40\% 时，A、C 两型均无法完成该服务区的安全往返，B 型也只能携带约 $1.397\,\mathrm{kg}$。该现象解释了高安全余量下架次数迅速增加直至整体不可行的原因。
\begin{figure}[H]
  \centering
  \includegraphics[width=0.85\textwidth]{figures/s800.png}
  \caption{S008 安全载荷敏感性}
  \label{fig:q1-s008-sensitivity}
\end{figure}

\subsubsection{求解正确性与精确性检验}

为避免仅依赖单一算法得到的结果，本文对问题一进行了两类独立核验。首先，基于 DEM 对节点间航段进行逐像元地形检查，确保巡航海拔与地形最高点的计算一致；其次，以动态规划枚举得到的可行组批模式为基础，对每个服务区单独建立 0-1 整数规划模型，并检查最少架次数与对应最小能耗。15 个服务区的整数规划最优间隙均为 0，且结果全部与动态规划一致。

\begin{table}[H]
  \centering
  \small
  \setlength{\tabcolsep}{5pt}
  \renewcommand{\arraystretch}{1.15}
  \caption{动态规划与整数规划交叉验证结果}
  \label{tab:q1-validation}
  \begin{tabular*}{0.92\linewidth}{@{\extracolsep{\fill}}crrcc@{}}
    \toprule
    服务区 & 最少架次数 & \shortstack{对应能耗\\ /$\mathrm{kWh}$} & MILP gap & 一致性 \\
    \midrule
    S001 & 2 & 5.9134 & 0.0 & 通过 \\
    S002 & 2 & 8.4346 & 0.0 & 通过 \\
    S003 & 2 & 8.5449 & 0.0 & 通过 \\
    S004 & 1 & 6.1529 & 0.0 & 通过 \\
    S005 & 1 & 4.2222 & 0.0 & 通过 \\
    S006 & 1 & 2.5980 & 0.0 & 通过 \\
    S007 & 1 & 3.4109 & 0.0 & 通过 \\
    S008 & 1 & 5.8360 & 0.0 & 通过 \\
    S009 & 1 & 2.0962 & 0.0 & 通过 \\
    S010 & 1 & 1.8232 & 0.0 & 通过 \\
    S011 & 1 & 1.0740 & 0.0 & 通过 \\
    S012 & 1 & 2.7524 & 0.0 & 通过 \\
    S013 & 1 & 1.8248 & 0.0 & 通过 \\
    S014 & 1 & 2.1404 & 0.0 & 通过 \\
    S015 & 1 & 2.3075 & 0.0 & 通过 \\
    \bottomrule
  \end{tabular*}
\end{table}

从算法规模看，单个服务区枚举的可行模式数最多为 201 个，记忆化动态规划实际访问的状态数最多为 165 个，因此问题一可以在很短时间内精确求解。程序记录的本次问题一总运行时间约为 $0.679\,\mathrm{s}$。该结果也说明，利用“同属性货箱归类 + 状态压缩”的方式能够显著减少离散组合规模。

\subsection{问题一求解结果总结与讨论}

针对单服务区直接往返运输能力与不可拆货箱组批问题，本文在不考虑实体无人机数量及共享电池调度的条件下，依据题设地形净空、飞行时间与运输能耗计算规则，建立最大安全载荷模型，并综合货箱质量、装载体积和返航安全能量约束，采用动态规划求解各服务区的货箱组批方案。针对问题一的各项要求，主要结果如下。

\noindent\textbf{（1）三种机型在不同服务区的最大安全载荷。}

在默认 $20\%$ 返航安全余量下，A 型无人机在全部 15 个服务区的最大安全载荷均为 $25\,\mathrm{kg}$；B 型在 S008 的最大安全载荷为 $28.801\,\mathrm{kg}$，其余服务区均为 $30\,\mathrm{kg}$；C 型在 S002、S003、S004、S008 和 S012 的最大安全载荷分别为 $68.860$、$68.206$、$63.697$、$58.904$ 和 $68.117\,\mathrm{kg}$，其余服务区均为 $80\,\mathrm{kg}$。

上述结果表明，山区单点往返任务的最大安全载荷不仅取决于机型额定载质量，还受到水平航程、地形引起的爬升需求及返航安全余量的共同约束。在当前参数下，A 型均可达到额定载荷上限，而 B、C 两型在部分服务区受到能量约束。其中，S008 同时使 B、C 两型的最大安全载荷低于额定上限，是本算例中具有代表性的受限服务区。因此，不能直接以额定载质量替代各服务区的实际安全载荷进行货箱组批。

\noindent\textbf{（2）满足货箱与运输约束的组批方案。}

在默认条件下，80 个货箱被划分为 18 个直接往返架次，每个货箱恰好安排一次，各架次均满足载质量、装载体积及返航安全能量约束。其中，B 型、C 型各承担 9 个架次，A 型未被采用。具体而言，S001 由两个 C 型架次完成，载荷分别为 $76\,\mathrm{kg}$ 和 $78\,\mathrm{kg}$；S002、S003 均由一个 B 型架次和一个 C 型架次完成，载荷分别为 $25\,\mathrm{kg}$ 和 $56\,\mathrm{kg}$；S004 至 S008 各由一个 C 型架次完成；S009 至 S015 各由一个 B 型架次完成。各架次的货箱组成、载荷、体积、能耗及返航剩余电量见前文组批结果表。

A 型未被选中并不意味着其无法完成安全往返，而是质量、体积约束与优化目标共同作用的结果。例如，由饮用水、医疗物资和应急食品构成的典型组合总质量为 $25\,\mathrm{kg}$，但总体积为 $0.067\,\mathrm{m}^{3}$，超过 A 型 $0.060\,\mathrm{m}^{3}$ 的装载体积，却未超过 B 型 $0.073\,\mathrm{m}^{3}$ 的装载体积。因此，即使货箱组合满足质量上限，也未必能够由相应机型单架次完成运输，组批时必须同时检查质量与体积约束。

\noindent\textbf{（3）架次数、运输能耗与累计作业时间的优化结果及权衡。}

本文采用“首先最小化架次数，其次最小化总运输能耗，最后最小化累计作业时间”的字典序优先关系。所得方案共需 18 个架次，总运输能耗为 $59.1313\,\mathrm{kWh}$，累计作业时间为 $9.104\,\mathrm{h}$。这里的累计作业时间是全部架次作业时长之和，并非考虑多架无人机并行执行后的任务完工时间。

对照实验表明，将能耗改为第一优先目标后，总运输能耗降至 $59.0339\,\mathrm{kWh}$，降幅约为 $0.165\%$；但架次数增加至 19，增幅约为 $5.56\%$，累计作业时间增加至 $9.567\,\mathrm{h}$，增幅约为 $5.08\%$。以累计作业时间为第一目标时，所得方案与本文采用的字典序方案一致。因此，在本算例中，架次数优先能够以较小的能耗代价减少运输批次数和累计作业时间。该结果说明本文方案是在明确优先关系下取得的折中解，而不是三个指标分别达到各自独立最优的方案。

\noindent\textbf{（4）返航安全余量变化对最大安全载荷与组批结果的影响。}

随着返航安全余量提高，可用于运输的能量预算减少。对于固定的机型与服务区，最大安全载荷随余量增加不会上升，而是可能保持额定上限、逐步下降，直至该组合无法满足空载安全往返要求。以 S008 为例，安全余量由 $20\%$ 提高至 $30\%$ 时，A、B、C 三型的最大安全载荷分别由 $25.000$、$28.801$、$58.904\,\mathrm{kg}$ 降至 $14.213$、$21.001$、$36.814\,\mathrm{kg}$；余量达到 $35\%$ 时，A 型已无法满足空载安全往返要求，C 型的最大安全载荷降至 $15.683\,\mathrm{kg}$。

在组批层面，已测试的 $10\%$、$15\%$ 和 $20\%$ 余量下，最优方案均保持 18 个架次，总运输能耗均为 $59.1313\,\mathrm{kWh}$。余量提高至 $25\%$、$30\%$ 和 $35\%$ 时，架次数分别增加至 19、20 和 25，总运输能耗分别增加至 $61.073$、$67.228$ 和 $75.608\,\mathrm{kWh}$，累计作业时间分别增加至 $9.602$、$10.163$ 和 $12.607\,\mathrm{h}$。当余量为 $40\%$ 时，在现有机型与单点往返规则下，已不存在能够完成全部货箱交付的可行组批方案。

综上，问题一给出了三种机型对 15 个服务区的最大安全载荷、满足约束的完整货箱组批方案、不同优化优先级下的指标权衡，以及返航安全余量对载荷能力和组批结果的影响。
\FloatBarrier

\subsection{算法示例}

数学建模求解算法示例：
\begin{center}
\begin{minipage}{0.8\textwidth}
\begin{algorithm}[H]%[!htp]
\caption{算法的名字} %算法的名字
{\bf 输入:} %算法的输入， \hspace*{0.02in}用来控制位置，同时利用 \\ 进行换行
input parameters A, B, C\\
{\bf 输出:} %算法的结果输出
output result
\begin{algorithmic}[1]
\State some description 算法介绍 % \State 后写一般语句
\For{condition} % For 语句，需要和EndFor对应
  \State ...
  \If{condition} % If 语句，需要和EndIf对应
    \State ...
  \Else
    \State ...
  \EndIf
\EndFor
\While{condition} % While语句，需要和EndWhile对应
  \State ...
\EndWhile
\State \Return result
\end{algorithmic}
\end{algorithm}
\end{minipage}
\end{center}
\vspace{2ex}

\subsection{表格}

使用 tabular 环境, 如下表格: 表~\ref{tab:foo}. 通过 \verb|autoref| 引用表格: \autoref{tab:foo}.

这里可使用命令 \verb|\tabcolsep| 和 \verb|\arraystretch| 分别控制列间距和行间距.

\begin{table}[htp!]
\centering
\setlength{\tabcolsep}{12pt}  % 6pt standard
\renewcommand{\arraystretch}{1.2}
\caption{学术活动安排样例}
\label{tab:foo}
\begin{tabular}{|c|c|c|c|}
\hline
\textbf{日期}  & \textbf{地点} & \textbf{活动名称} & \textbf{备注} \\ \hline
2024年8月1日      & 上海       & 学术研讨会      & 主题：人工智能 \\ \hline
2024年8月15日    & 北京       & 学术交流会      & 重点：数学建模 \\ \hline
2024年9月1日      & 深圳       & 研究研讨会      & 主题：数据科学 \\ \hline
2024年10月15日  & 广州       & 创新论坛         & 重点：科技创新 \\ \hline
\end{tabular}
\end{table}


\clearpage

三线表
\begin{table}[htp!]
\newcolumntype{L}{X}
\newcolumntype{C}{>{\centering \arraybackslash}X}
\newcolumntype{R}{>{\raggedright \arraybackslash}X}
\centering
\caption{某校学生升高体重样本}
\label{tab:heightweight}
\begin{tabularx}{0.9\textwidth}{CCCC}
\toprule
序号 & 年龄 & 身高 & 体重 \\
\midrule
001 & 15 & 156 & 42 \\
002 & 16 & 158 & 45 \\
003 & 14 & 162 & 48 \\
004 & 15 & 163 & 50 \\
\cmidrule{2-4}
平均 & 15 & 159.75 & 46.25 \\
\bottomrule
\end{tabularx}
\end{table}

% 某行业产量与生产费用的数据
% \begin{table}[htp!]
% \centering
% \caption{某行业产量与生产费用的数据}%\label{}
% \newcolumntype{Y}{>{\centering\arraybackslash}X}
% \newcolumntype{Z}{!{\vline}@{\color{white}\vrule width \doublerulesep}!{\vrule}}%自定义列格式(双线)
% \begin{tabularx}{0.94\textwidth}{c|c|YZc|c|Y}
% \Xhline{0.9pt}
% 企业编号&  产量(台)&生产费用(万元)&企业编号&产量(台)&生产费用(万元)\\\Xcline{1-3}{0.6pt}\Xcline{4-6}{0.6pt}
% 1 & 40 & 130 & 7&  84&  165\\
% 2 & 42 & 150 & 8&  100&  170\\
% 3 & 50 & 155 & 9&  116&  167\\
% 4 & 55 & 140 & 10&  125&  180\\
% 5 & 65 & 150 & 11&  130&  175\\
% 6 & 78 & 154 & 12&  140&  185\\\Xhline{0.72pt}
% \end{tabularx}
% \end{table}

研究生数学建模2019年F题结果示例
\begin{table}[htp!]
\centering
\caption{问题1结果1 (左) 与 问题2结果 (右)}
\begin{minipage}[t]{0.48\linewidth}
\renewcommand\arraystretch{1.2} %定义表格高度
\newcolumntype{Y}{>{\centering\arraybackslash}X}
\begin{tabularx}{0.9\textwidth}{|Y|Y|Y|}
  \hline
  数据集1  &  数据集1 & 数据集2  \\
 \hline
  A问题1   & A问题1 &  A问题1     \\
  \hline
  503   &  503     & 163     \\
  294   &  200     & 114      \\
  91    &  80      & 8     \\
  607   &  237     & 309      \\
  540   &  170     & 305    \\
  250   &  278     & 123    \\
  340   &  369     & 45      \\
  277   &  214     & 160    \\
  B     &  397     & 92    \\
        &  B       & 93    \\
        &          & 61        \\
        &          & 292       \\
        &          & B         \\
104861  & 103518   & 109342 \\
\hline
\end{tabularx}
\end{minipage}
\begin{minipage}[t]{0.48\linewidth}
\renewcommand\arraystretch{1.2} %定义表格高度
\newcolumntype{Y}{>{\centering\arraybackslash}X}
\begin{tabularx}{0.9\textwidth}{|Y|Y|Y|}
  \hline
  数据集1  &  数据集1 & 数据集2  \\
  \hline
  A问题2  &  A问题2  & A问题2    \\
  \hline
   503    & 503     & 163  \\
  294    & 200     & 114   \\
   91     & 80     & 8   \\
   607    & 237     & 309   \\
   540    & 170    & 305  \\
   250    & 278    & 123  \\
   340    & 369     & 45   \\
   277    & 214    & 160  \\
  B      & 397   & 92 \\
        &  B     &  93      \\
         &         &  61      \\
        &         &   292     \\
        &         &   B     \\
 104917  &  103563  &109427 \\
\hline
\end{tabularx}
\end{minipage}
\end{table}

\clearpage
\begin{table}[htp!]
%\small
\centering
\renewcommand\arraystretch{1.2} %定义表格高度
\newcolumntype{Y}{>{\centering\arraybackslash}X}
\caption{问题3结果}
\begin{tabularx}{0.9\textwidth}{|Y|Y|Y|Y|Y|Y|}
  \hline %定义表格宽度
  数据集1  &  数据集1  &  \multicolumn{2}{c|}{数据集2 (无问题点)}  & \multicolumn{2}{c|}{数据集2 (有问题点)} \\
  \hline
  A问题3   & A问题3   &  \multicolumn{2}{c|}{A问题3}    &  \multicolumn{2}{c|}{A问题3}      \\
  \hline
  503      & 503     & 169      & 73      & 169     & 73   \\
  \hline
  69       & 69      & 322      & 249     & 322     & 249   \\
  \hline
  506      & 506      & 270     & 274     & 270     & 274   \\
  \hline
  371      & 371      & 89      & 12      & 89      & 12   \\
  \hline
  183      & 183      & 236     & 216     & 236     & 216  \\
  \hline
  194      & 194      & 132     & 16      & 132     & 16   \\
  \hline
  450      & 450      & 53      & 282     & 53      & 282   \\
  \hline
  286      & 113      & 112     & 84      & 112     & 141  \\
  \hline
  485      & 485      &  268    & 287     & 268     & 291 \\
  \hline
 \red{B (9D)}~~     & 248      & 250     & 99      & 250     &161 \\
 \hline
   & \red{B (10D)}   & 243     & \red{B (21D)}   & 243    & \red{B (21D)} \\
   \hline
    &          &         &         &         &       \\
    \hline
  104861m  & 103518m  &         & 168924m  &         &161650m \\
\hline
\end{tabularx}
\end{table}


\clearpage

\section{问题二的建模与求解}

\subsection{问题二分析}
问题二在问题一的基础上，由单点往返扩展为\textbf{异构无人机多点、多架次联合调度}。其主要难点在于货箱组批、服务区访问顺序、无人机选择、架次开始时间以及共享电池周转相互耦合，若直接对所有决策变量同时搜索，将产生较大的组合规模。

因此，本问可采用\textbf{“可行架次生成 + 资源时序调度”}的分层求解思路。首先依据问题一得到的载荷与能耗计算方法，生成满足重量、体积、飞行能量和返航安全要求的候选运输架次，并确定每个架次可访问的服务区及其访问顺序；随后将这些候选架次分配给具体无人机和共享电池，并确定各架次开始时刻。调度过程中重点处理医疗物资、首批保障物资等具有较强时效要求的货箱，同时保证无人机任务占用、电池充电及再次使用之间不存在时间冲突。最终以\textbf{配送及时性、任务完成时间、总能耗和架次数}作为主要评价指标，形成兼顾救援效率与资源利用率的运输调度方案。


\section{问题三的建模与求解}
\subsection{问题三分析}
问题三的关键是将问题二中的运输调度进一步与\textbf{连续通信保障}相结合。山区环境中，运输无人机与固定网关之间的通信状态会随位置、高度和地形遮挡不断变化，因此不能仅判断服务区处是否具有通信条件，而应对无人机的整个飞行轨迹进行通信可用性检测。

首先根据运输无人机在爬升、巡航、下降和投送阶段的空间轨迹，结合DEM地形进行视距遮挡判断，并利用链路预算计算无人机与固定网关之间的通信状态。对于无法直连的轨迹区段，引入中继无人机提供通信保障。由于中继悬停位置属于连续空间变量，直接联合优化难度较大，因此可在满足地形和高度限制的区域内构造\textbf{候选中继位置集合}，预先计算各候选位置对不同运输轨迹和时间区段的通信覆盖关系，将连续空间问题转化为有限候选点上的时空覆盖问题。

在此基础上，将运输架次、中继位置、中继服务时段和能源组件进行联合调度，保证每一时刻运输无人机至少能够通过固定网关直连或中继链路保持通信。故本问的核心突破口在于\textbf{先建立轨迹级通信可行性矩阵，再进行运输--中继协同调度}，从而在保证通信连续性的同时，兼顾任务完成时间、运输与中继能耗以及设备使用规模。

\section{问题四的建模与求解}
\subsection{问题四分析}
问题四是在问题三既定联合调度方案的基础上研究\textbf{任务分区与资源独立配置}问题。由于题目要求保持原有货箱组批、访问顺序、运输任务和通信保障关系不变，因此本问并不重新规划运输路线，而是需要识别哪些服务区之间因共享同一运输架次或通信任务而具有不可分割的关联关系。

为此，可首先构建\textbf{服务区任务关联图}：将各服务区视为节点，若多个服务区出现在同一运输架次中，则在其之间建立关联，并将相互关联的服务区合并为不可拆分的任务单元。在此基础上，分别对这些任务单元进行两分区和三分区，使各组在任务量和资源需求上尽可能均衡。

完成分区后，需要根据每个任务组中的运输与中继任务时间区间，计算运输无人机、共享电池、中继无人机和能源组件的\textbf{最大并发占用量}，而不能简单累加任务总数，以此确定各分区独立运行时真正需要配置的最小资源数量。最终从资源配置规模、资源冗余、工作量均衡程度及现有库存缺口等方面比较两种分区方案，并识别导致资源不足的关键任务和时间段。

\section{模型评价}
\subsection{模型的优缺点}
\subsection{模型的改进方向}


%参考文献   手工录入
%\begin{thebibliography}{9}%宽度9
% \bibitem{bib:one} ....
% \bibitem{bib:two} ....
%\end{thebibliography}

%采用bibtex方案
\cite{mittelbach_latex_2004,wright_latex3_2009,beeton_unicode_2008,vieth_experiences_2009}

\bibliographystyle{gmcm}
\bibliography{reference}


\clearpage
%附录
\begin{appendices}
%\setcounter{page}{1} %如果需要可以自行重置页码。
\section{MATLAB 源程序}
\renewcommand{\thesubsection}{A\thinspace.\thinspace\arabic{subsection}}
\subsection{第1问程序}
\vspace{-2ex}

\begin{Matlab}{code.m}
clear all
kk=2;
[mdd,ndd]=size(dd);
while ~isempty(V)
    [tmpd,j]=min(W(i,V));
    tmpj=V(j);
    for k=2:ndd
        [tmp1,jj]=min(dd(1,k)+W(dd(2,k),V));
        tmp2=V(jj);
        tt(k-1,:)=[tmp1,tmp2,jj];
    end
    tmp=[tmpd,tmpj,j;tt];
    [tmp3,tmp4]=min(tmp(:,1));
    if tmp3==tmpd,
        ss(1:2,kk)=[i;tmp(tmp4,2)];
    else
        tmp5=find(ss(:,tmp4)~=0);
        tmp6=length(tmp5);
        if dd(2,tmp4)==ss(tmp6,tmp4)
            ss(1:tmp6+1,kk)=[ss(tmp5,tmp4);tmp(tmp4,2)];
        else, ss(1:3,kk)=[i;dd(2,tmp4);tmp(tmp4,2)];
        end
    end
    dd=[dd,[tmp3;tmp(tmp4,2)]];
    V(tmp(tmp4,3))=[];
    [mdd,ndd]=size(dd);kk=kk+1;
end;
S=ss; D=dd(1,:);
\end{Matlab}
\vspace{2ex}

\clearpage
\section{Python 源程序}
\renewcommand{\thesubsection}{B\thinspace.\thinspace\arabic{subsection}}
\subsection{第2问程序}
\vspace{-2ex}
\begin{Python}{mip1.py}
# This example formulates and solves the following simple MIP model:
#  maximize
#        x +   y + 2 z
#  subject to
#        x + 2 y + 3 z <= 4
#        x +   y       >= 1
#        x, y, z binary

# import gurobipy as gp
from gurobipy import * #GRB
try:
    # Create a new model
    m = Model("mip1")
    # Create variables
    x = m.addVar(vtype=GRB.BINARY, name="x")
    y = m.addVar(vtype=GRB.BINARY, name="y")
    z = m.addVar(vtype=GRB.BINARY, name="z")
    # Set objective
    m.setObjective(x + y + 2 * z, GRB.MAXIMIZE)
    # Add constraint: x + 2 y + 3 z <= 4
    m.addConstr(x + 2 * y + 3 * z <= 4, "c0")
    # Add constraint: x + y >= 1
    m.addConstr(x + y >= 1, "c1")
    # Optimize model
    m.optimize()
    for v in m.getVars():
        print('%s %g' % (v.varName, v.x))
    print('Obj: %g' % m.objVal)

except GurobiError as e:
    print('Error code ' + str(e.errno) + ': ' + str(e))

except AttributeError:
    print('Encountered an attribute error')
\end{Python}

\end{appendices}


\end{document}

#        x +   y       >= 1
#        x, y, z binary

# import gurobipy as gp
from gurobipy import * #GRB
try:
    # Create a new model
    m = Model("mip1")
    # Create variables
    x = m.addVar(vtype=GRB.BINARY, name="x")
    y = m.addVar(vtype=GRB.BINARY, name="y")
    z = m.addVar(vtype=GRB.BINARY, name="z")
    # Set objective
    m.setObjective(x + y + 2 * z, GRB.MAXIMIZE)
    # Add constraint: x + 2 y + 3 z <= 4
    m.addConstr(x + 2 * y + 3 * z <= 4, "c0")
    # Add constraint: x + y >= 1
    m.addConstr(x + y >= 1, "c1")
    # Optimize model
    m.optimize()
    for v in m.getVars():
        print('%s %g' % (v.varName, v.x))
    print('Obj: %g' % m.objVal)

except GurobiError as e:
    print('Error code ' + str(e.errno) + ': ' + str(e))

except AttributeError:
    print('Encountered an attribute error')
\end{Python}

\end{appendices}


\end{document}
